% WAFC -- manuscrito, versao k = 5 (E5h, 2026-10-06; copia de ms_4.tex, que fica
% intacta; a rodada nao foi aceita, entao a marcacao de k = 2 a k = 5 segue em colR1).
% Template da Statistica Sinica (manuscript/ss-template/SS-template-bib.tex), decisao D5.
% Secoes 1 a 6; as Secoes 5 (Simulation) e 6 (Application) entraram na E5h sem os
% numeros, com marcadores [E4.3], [E4.4] e [E6.2] no lugar deles; a tabela e as figuras
% vem por \input de tables/ e figures/ (D21). A Secao 7 (Discussion) e da E5b.
% As provas estao em supp_5.tex.
% (Cabecalho da k = 4, mantido: versao k = 4 (E5d; copia de k = 3, alteracoes marcadas
% em colR1). Secoes 1 a 4; as Secoes 5 a 7 entram em E5b.)
% Notacao: docs/notacao.md (congelada em E1.1; a penalidade em blocos no seu S10).
% Os enunciados vem de derivations/01 a 08; o comentario acima de cada um diz de onde.
% Em k = 4 (D44 a D51) o estimador passa a ser o block lasso balanceado, e a teoria
% do lasso coordenado sai do corpo para a Secao S8 do supp (D51(a)). Enunciado,
% prova ou matematica deslocada que sai inteiro fica no lugar, em cinza e sem
% numero, aberto por "[removed in k = 4]" (D51(b)); o novo vem logo depois.
%
% Mapa da numeracao global (docs/TAREFA.md, S5) para os rotulos deste arquivo, k = 4:
%   Proposicao 1            -> Proposition 1  (prop:identif)
%   Lema 1                  -> Lemma 1        (lem:basis)
%   Lema 2 e Corolario 1    -> Lemma 2        (lem:bias)
%   Proposicao 3            -> Proposition 2  (prop:gram)
%   Corolario 11            -> Theorem 1      (thm:main; D47(a), emenda D16)
%   Teorema 3               -> Theorem 2      (thm:oracle); o Lema 14 e o Lemma S5.2
%                              do supp (lem:calib-blocks), a Observacao 2 (variante
%                              branca) e a Remark S5.2, o Corolario 9 e o
%                              Corollary S5.1, a Proposicao 7 e a Proposition S5.1
%   Teorema 4               -> Theorem 3      (thm:sieve)
%   Corolario 10            -> Corollary 1    (cor:components)
%   Corolario 14            -> Proposition 3  (prop:slow; D47(c), D49(b)); os itens
%                              (iii) e (iv) de E1.14 sao a Proposition S6.1 do supp
%                              (prop:slow-ceiling; E5e, D52)
%   Corolario 15            -> Proposition S6.2 do supp (prop:truncated; E5f, D54),
%                              citada numa frase depois da Proposition 3
%   Teorema 5 e Lema 17     -> Theorem S6.1 (thm:lower) e Lemma S6.3 (lem:hypercube)
%                              do supp (E5g, D57), no lugar da Remark S6.1
%                              (rem:lower, em cinza); citado na terceira leitura
%                              do Theorem 1 e na introducao. O Corolario 16 nao
%                              entra como enunciado: e a leitura depois do
%                              Theorem S6.1, com a uniformidade lida na prova do
%                              Theorem 1 no supp
%   Lema 15                -> Lemma 3        (lem:weak; o rotulo ficou com o nome antigo)
%   Corolario 12, Hipotese S com rho_n^G
%                           -> Corollary 2 (cor:threshold), Assumption 6 (ass:sep);
%                              o Lema 13 e o Lemma S7.1 do supp (lem:sandwich)
%   Corolario 13            -> Corollary 3    (cor:thrisk; D45); o Lema 16 e o
%                              Lemma S7.2 do supp (lem:thrisk)
%   Teoria do lasso coordenado (Corolario 5, Teorema 1, Teorema 2, Corolario 4,
%   Proposicao 4 emendada por D49 e com os itens (iii) e (iv) de E1.14 (E5e),
%   Lema 9, Corolario 8): Secao S8 do supp,
%   como Theorems S8.1 a S8.3, Corollaries S8.1 e S8.2, Proposition S8.1 e
%   Lemma S8.1 (D51(a)). A Proposicao 8 (E1.15) e a Proposition S8.2
%   (prop:lasso-truncated; E5f, D54).
%   Os rotulos sec:sieve, thm:sieve e lem:weak ficaram com o nome antigo (D33 muda
%   o texto, nao os rotulos; em k = 4 o conteudo de thm:sieve e lem:weak mudou).
\documentclass[12pt]{article} %***
\usepackage[sectionbib]{natbib}
\usepackage{array,epsfig,fancyheadings,rotating}
\usepackage[]{hyperref}
%%%%%%%%%%%%%%%%%%%%%%%%%%%%%%%%%%%%
\usepackage{sectsty, secdot}
\sectionfont{\fontsize{12}{14pt plus.8pt minus .6pt}\selectfont}
\renewcommand{\theequation}{\thesection\arabic{equation}}
\subsectionfont{\fontsize{12}{14pt plus.8pt minus .6pt}\selectfont}
%%%%%%%%%%%%%%%%%%%%%%%%%%%%%%%%%%%%%%%%%%%%%%%%%%%%%%%%%%%%%%%%%%%%%%%%%%%%%%%%%%%%%%%

\textwidth=31.9pc
\textheight=46.5pc
\oddsidemargin=1pc
\evensidemargin=1pc
\headsep=15pt
\topmargin=.6cm
\parindent=1.7pc
\parskip=0pt

\usepackage{amsmath}
\usepackage{amssymb}
\usepackage{amsfonts}
\usepackage{multirow}
\usepackage{amsthm}

% Marcacao de alteracoes a partir de k = 2 (docs/instrucoes.md, S3): a cor da
% rodada corrente e colR1, e o ulem fornece o \sout das remocoes sugeridas.
% A versao k = 1 nao leva marcacao, porque nao ha versao anterior.
\usepackage{xcolor}
\usepackage[normalem]{ulem}
\definecolor{colR1}{RGB}{0,90,200}

\setcounter{page}{1}
\newtheorem{theorem}{Theorem}
\newtheorem{lemma}{Lemma}
\newtheorem{corollary}{Corollary}
\newtheorem{proposition}{Proposition}
\theoremstyle{definition}
\newtheorem{definition}{Definition}
\newtheorem{example}{Example}
\newtheorem{remark}{Remark}
\newtheorem{assumption}{Assumption}
% Enunciados removidos em k = 4 (D51(b)): os mesmos ambientes, sem numero, para
% que o texto antigo fique no lugar sem deslocar a numeracao do novo.
\newtheorem*{assumptionold}{Assumption}
\theoremstyle{plain}
\newtheorem*{theoremold}{Theorem}
\newtheorem*{lemmaold}{Lemma}
\newtheorem*{corollaryold}{Corollary}
\newtheorem*{propositionold}{Proposition}
\pagestyle{fancy}

% --- macros -------------------------------------------------------------
% Implementam docs/notacao.md, como derivations/macros.tex, com tres
% diferencas exigidas pelo estilo da revista (S5 de ss-instrucoes-autores.md):
% E(.) e P(.) em romano e com parenteses, no lugar do \mathbb{E}[.] das
% derivacoes, e B_X no lugar dos dois simbolos C_X (E1.3) e B_X (E1.4) para a
% cota das covariaveis lineares.
\newcommand{\R}{\mathbb{R}}
\newcommand{\E}{\mathrm{E}}
\newcommand{\Prob}{\mathrm{P}}
\newcommand{\Var}{\mathrm{Var}}
\newcommand{\T}{^{\top}}
\newcommand{\norm}[1]{\left\lVert #1 \right\rVert}
\newcommand{\normn}[1]{\left\lVert #1 \right\rVert_{n}}
\newcommand{\normL}[1]{\left\lVert #1 \right\rVert_{L_2(P)}}
\newcommand{\inner}[2]{\langle #1, #2 \rangle}
\newcommand{\cv}{\ell}
\newcommand{\md}{m}
\newcommand{\lev}{j}
\newcommand{\tsl}{k}
\newcommand{\jzero}{j_{0}}
\newcommand{\NJ}{N_{J}}
\newcommand{\bX}{\mathbf{X}}
\newcommand{\bU}{\mathbf{U}}
\newcommand{\bZ}{\mathbf{Z}}
\newcommand{\fcoef}[1]{\beta_{#1}}
\newcommand{\lvl}[1]{c_{#1}}
\newcommand{\gcomp}[2]{g_{#1#2}}
\newcommand{\wav}[2]{\psi_{#1#2}}
\newcommand{\VJ}{V_{J}}
\newcommand{\WJ}{W_{J}}
\newcommand{\PiJ}{\Pi_{J}}
\newcommand{\coef}[4]{\theta_{#1#2,#3#4}}
\newcommand{\coefs}[4]{\theta^{*}_{#1#2,#3#4}}
\newcommand{\btheta}{\boldsymbol{\theta}}
\newcommand{\thetastar}{\btheta^{*}}
\newcommand{\thetahat}{\widehat{\btheta}}
\newcommand{\bbeta}{\boldsymbol{\beta}}
\newcommand{\betastar}{\bbeta^{*}}
\newcommand{\bc}{\mathbf{c}}
\newcommand{\chat}{\widehat{\mathbf{c}}}
\newcommand{\Gram}{\Sigma}
\newcommand{\Gramhat}{\widehat{\Sigma}}
\newcommand{\pen}{\lambda}
\newcommand{\penn}{\lambda_{n}}
\newcommand{\Szero}{S_{0}}
\newcommand{\Besov}[3]{B^{#1}_{#2,#3}}
\newcommand{\seff}{s'}
\newcommand{\Amat}{\mathbf{A}}
\newcommand{\Bmat}{\mathbf{B}}
\newcommand{\Btil}{\widetilde{\Bmat}}
\newcommand{\PA}{\Pi_{\Amat}}
\newcommand{\MA}{M_{\Amat}}
\newcommand{\bveps}{\boldsymbol{\varepsilon}}
\newcommand{\bb}{\mathbf{b}}
\newcommand{\bY}{\mathbf{Y}}
\newcommand{\bPsi}{\boldsymbol{\Psi}}
\newcommand{\smax}{\widehat{\sigma}_{\max}}
\newcommand{\gtil}{\widetilde{\gamma}}
\newcommand{\Mstar}{M^{\star}}
\newcommand{\lmin}{\lambda_{\min}}
\newcommand{\lmaxx}{\lambda_{\max}}
\newcommand{\opnorm}[1]{\left\lVert #1 \right\rVert_{\mathrm{op}}}
% --- a penalidade em blocos (docs/notacao.md, S10; copiadas das macros locais
% de derivations/08-blocos.tex e 06-selecao-limiar.tex em k = 4) ---------------
\newcommand{\thbar}{\bar{\btheta}}
\newcommand{\bbar}{\bar{\bb}}
\newcommand{\vbar}{\bar{v}}
\newcommand{\Gcal}{\mathcal{G}}
\newcommand{\Gzero}{\mathcal{G}_{0}}
\newcommand{\Gbar}{\bar{\mathcal{G}}}
\newcommand{\Hcal}{\mathcal{H}}
\newcommand{\Wcal}{\mathcal{W}}
\newcommand{\bn}{b_{n}}
\newcommand{\jst}{j^{*}}
\newcommand{\nGw}[1]{\lVert #1\rVert_{\mathcal{G},w}}
\newcommand{\nGone}[1]{\lVert #1\rVert_{\mathcal{G},1}}
\newcommand{\Psitil}{\widetilde{\Psi}}
\newcommand{\Rid}{R_{\mathcal{G}}}
\newcommand{\lamG}{\lambda^{\mathcal{G}}_{n}}
\newcommand{\lamGone}{\lambda^{\mathcal{G}}_{n,1}}
\newcommand{\lamGp}{\lambda^{\mathcal{G},+}_{n}}   % E1.14 (E5e): \lamG inflado acima do teto
\newcommand{\lamz}[1]{\lambda_{0,#1}}
\newcommand{\lamzbar}[1]{\bar{\lambda}_{0,#1}}
\newcommand{\rhoG}{\rho^{\mathcal{G}}_{n}}
\newcommand{\rhoGt}{\tilde{\rho}^{\mathcal{G}}_{n}}
\newcommand{\Tw}{\mathcal{T}_{\mathcal{G},w}}
\newcommand{\AG}{A^{\mathcal{G}}_{s,\pi}}
% grupos de marcacao de D51(b): o texto antigo, inteiro, em cinza e sem tachado
\newcommand{\removedk}{\par\noindent[removed in k = 4]\par}

%%%%%%%%%%%%%%%%%%%%%%%%%%%%%%%%%%%%%%%%%%%%%%%%%%%%%%%%%%%%%%%%%%%%%%%%%%%%%%
\pagestyle{fancy}
\def\n{\noindent}
\lhead[\fancyplain{} \leftmark]{}
\chead[]{}
\rhead[]{\fancyplain{}\rightmark}
\cfoot{}

%%%%%%%%%%%%%%%%%%%%%%%%%%%%%%%%%%%%%%%%%%%%%%%%%%%%%%%%%%%%%%%%%%%%%%%%%%%%%%
\begin{document}

\renewcommand{\baselinestretch}{2}

\markright{ \hbox{\footnotesize\rm Statistica Sinica
  }\hfill\\[-13pt]
  \hbox{\footnotesize\rm
  }\hfill }

\markboth{\hfill{\footnotesize\rm AUTHOR(S)} \hfill}
{\hfill {\footnotesize\rm WAVELET ADDITIVE FUNCTIONAL COEFFICIENTS} \hfill}

\renewcommand{\thefootnote}{}
$\ $\par

%%%%%%%%%%%%%%%%%%%%%%%%%%%%%%%%%%%%%%%%%%%%%%%%%%%%%%%%%%%%%%%%%%%%%%%%%%%%%%

\fontsize{12}{14pt plus.8pt minus .6pt}\selectfont \vspace{0.8pc}
\centerline{\large\bf SPARSE WAVELET ESTIMATION OF}
\vspace{2pt}
\centerline{\large\bf ADDITIVE FUNCTIONAL COEFFICIENTS}
\vspace{.4cm}
\centerline{Author(s)}
\vspace{.4cm}
\centerline{\it Affiliation(s)}
\vspace{.55cm} \fontsize{9}{11.5pt plus.8pt minus.6pt}\selectfont

%%%%%%%%%%%%%%%%%%%%%%%%%%%%%%%%%%%%%%%%%%%%%%%%%%%%%%%%%%%%%%%%%%%%%%%%%%%%%%

\begin{quotation}
\noindent {\it Abstract:}
We study a varying coefficient model in which the effect of each linear
covariate is additive in several modulating covariates, and each additive
component is expanded in a periodized wavelet basis. The resulting design is
\textcolor{colR1}{an approximation space} of products between a linear covariate and a wavelet evaluated at a
modulator, and the components are estimated by a single \textcolor{gray}{\sout{lasso}}
\textcolor{colR1}{block lasso} problem in
which the levels of the coefficients are left unpenalized \textcolor{colR1}{and the
wavelet coefficients of each resolution level are grouped in chunks of order
$\log n$}. Three results are
established. The population Gram matrix of the product design has its
smallest eigenvalue bounded away from zero under a conditional second moment
condition on the linear covariates and a bounded joint density for the
modulators, with no independence between the two groups of covariates; the
estimator satisfies a nonasymptotic oracle inequality that needs no cone
condition, because the unpenalized levels are profiled out\textcolor{colR1}{, and
that holds for any weights of bounded ratio}; and, at a
resolution level chosen inside an explicit window, the prediction error and
the componentwise error attain the rate of nonlinear wavelet approximation\textcolor{gray}{\sout{ up
to a logarithmic factor}}\textcolor{colR1}{, with no logarithmic factor when the Besov
integrability index is at least two}\textcolor{colR1}{, which is minimax optimal}, under the same Besov condition that governs
\textcolor{colR1}{linear approximation}. The last result rests on \textcolor{gray}{\sout{the observation that the Besov
condition already implies compressibility of the}} \textcolor{gray}{\sout{wavelet}} \textcolor{gray}{\sout{coefficients}}\textcolor{colR1}{a
bound on the ideal risk over chunks that the Besov condition already implies},
so that adaptivity costs no additional assumption. \textcolor{colR1}{A threshold on the
norms of the fitted components recovers which modulators act on which
coefficients, and the thresholded estimator keeps the rate of the fit.} The
estimator is a single convex
program and is fit by \textcolor{gray}{\sout{coordinate descent on a sparse design matrix}}\textcolor{colR1}{standard
group lasso software}.
\textcolor{colR1}{An implementation in R and the code that reproduces every
numerical result of the paper are made available with it.}

\vspace{9pt}
\noindent {\it Key words and phrases:}
Additive coefficient model, Besov space, \textcolor{gray}{\sout{effective dimension,}} \textcolor{colR1}{group
lasso,} lasso, oracle
inequality, \textcolor{colR1}{thresholding,} varying coefficient model, wavelets.
\par
\end{quotation}\par

\def\thefigure{\arabic{figure}}
\def\thetable{\arabic{table}}

\renewcommand{\theequation}{\thesection.\arabic{equation}}
% o contador de equacoes do template nao reinicia por secao, o que produzia
% "(3.6)" logo depois de "(2.5)"; o \numberwithin reinicia.
\numberwithin{equation}{section}

\fontsize{12}{14pt plus.8pt minus .6pt}\selectfont

%%%%%%%%%%%%%%%%%%%%%%%%%%%%%%%%%%%%%%%%%%%%%%%%%%%%%%%%%%%%%%%%%%%%%%%%%%%%%%
\section{Introduction}
\label{sec:intro}

Applied problems often present an effect that is not constant across the
sample. The return to schooling changes with age and with experience, the
effect of temperature on electricity demand changes with the hour of the day
and with humidity, and the effect of a treatment changes with the age and the
body mass index of the patient. The varying coefficient model of
\citet{Hastie-Tibshirani-1993} and \citet*{Cai-Fan-Yao-2000} accommodates this
by letting the coefficient of each covariate be a function of a modulating
variable; see \citet{Fan-Zhang-2008} for a review. When more than one variable
modulates the effect, a fully nonparametric coefficient runs into the curse of
dimensionality, and the standard remedy is to make each coefficient additive
in the modulators, which is the additive coefficient model of
\citet{Xue-Yang-2006}, studied further by \citet{Xue-Liang-2010} and
\citet{Liu-Yang-2010}.

The estimators available for this model are built from polynomial splines and
from kernels, and are governed by a global smoothness index, so that a
component with a peak, a jump or a change of regime is either oversmoothed or
paid for everywhere in the domain. Wavelet bases combined with an $\ell_1$
penalty \textcolor{colR1}{\citep{tibshirani1996regression}} are the standard answer to this difficulty, because a function whose
smoothness is inhomogeneous is still economically represented by a few large
coefficients \citep{Donoho-Johnstone-1994, Donoho-Johnstone-1998,
Antoniadis-Fan-2001}. They have been used for varying coefficients without an
additive structure and without penalization \citep*{Zhou-You-2004,
Montoril-Morettin-Chiann-2018}, and for additive models without a multiplying
covariate \citep{Sardy-Tseng-2004, Zhang-Wong-2003, Haris-Simon-Shojaie-2018,
Amato-Antoniadis-DeFeis-Gijbels-2022, Sardy-Ma-2024}; penalized estimation in
the additive coefficient model itself has been carried out with grouped
penalties on spline bases \citep{Wei-Huang-Li-2011,
Antoniadis-Gijbels-LambertLacroix-2014, Ma-Carroll-Liang-Xu-2015}, again under
global smoothness.

\citet{Klopp-Pensky-2015} study the varying coefficient model in which a
single index modulates every coefficient. They expand each coefficient
function in an orthonormal basis, estimate the resulting array by a block
lasso, and obtain a nonasymptotic oracle inequality, a Besov rate that adapts
to inhomogeneous smoothness, and a matching minimax lower bound, under the
assumption that the covariates are independent of the index. The present paper
carries that \textcolor{gray}{\sout{programme}} \textcolor{colR1}{program} to the situation an applied problem usually presents,
in which several variables modulate a coefficient at once. We let each
coefficient be additive in the modulators, which keeps the estimator free of
the curse of dimensionality in their number, we allow the covariates to depend
on the modulators, and we \textcolor{gray}{\sout{penalize coefficient by coefficient rather than in
blocks, leaving the levels of the coefficients unpenalized. Each of these
costs something.}} \textcolor{colR1}{keep the block lasso but leave the levels of the
coefficients unpenalized and group the wavelet coefficients in chunks of order
$\log n$ within a resolution level, with weights whose ratio stays bounded, so
that the standardized weights of group lasso software are covered.}
Additivity produces cross terms between distinct modulators,
so the population Gram matrix no longer factorizes as a Kronecker product and
has to be bounded directly; dependence between the covariates and the
modulators replaces a marginal second moment by a conditional one; and the
unpenalized levels leave a term in the risk that \textcolor{colR1}{full} block penalization does not
incur. Our main result, Theorem~\ref{thm:main}, states that the estimator
attains the rate of \textcolor{gray}{\sout{nonlinear wavelet approximation up to a logarithmic
factor, under the same Besov assumption that governs}} \textcolor{gray}{\sout{linear approximation}}\textcolor{gray}{\sout{, and it
reduces to the rate of Klopp and Pensky when a single modulator is present.}}
\textcolor{colR1}{Klopp and Pensky, with no logarithmic factor when the Besov integrability
index is at least two, \textcolor{gray}{\sout{so that with a single modulator independent of the
covariates it is minimax optimal there}}
% E5g (D57): "where it is minimax optimal", pela cota inferior do Theorem S6.1.
where it is minimax optimal. It holds under weaker conditions than
theirs: the number of covariates is fixed, a regime their high dimensional
condition excludes, and the effective smoothness of the components need only
exceed $s/(2s+1)$, which is below one half, where theirs must exceed one half.
Components with effective smoothness one half, such as the inhomogeneous
functions of Section~\ref{sec:simulation}, are covered here and not there. \textcolor{gray}{\sout{A lower bound for
several modulators remains open.}}
% E5g (D57): a frase de D50 sobre a cota inferior, trocada (Theorem S6.1 do supp).
% D58 (pergunta 52(d)): "it" trocado por "the rate".
The lower bound that makes \textcolor{gray}{\sout{it}} the rate optimal is proved here for any number of
modulators and for covariates that depend on them; Klopp and Pensky have it
with a single modulator independent of the covariates.}
% E5h (k = 5): o "Section~5" acima passou a \ref{sec:simulation} (D50); o texto
% impresso nao muda. A Secao 5 diz que as funcoes nao homogeneas tem s' = 1/2.

We call the resulting procedure WAFC, for wavelet additive functional
coefficients. Its contributions are of \textcolor{gray}{\sout{three}} \textcolor{gray}{\sout{four}} \textcolor{colR1}{three} kinds. The first is the theory of
the additive product design: a design condition for the Gram matrix of the
products between a linear covariate and a wavelet of a modulator, which
carries a cross term between distinct modulators and does not assume
independence (Proposition~\ref{prop:gram}); an oracle inequality that needs no
cone condition, because the unpenalized levels are profiled out\textcolor{colR1}{, and
that holds for any weights of bounded ratio, including those of the software
that computes the estimator}
(Theorem~\ref{thm:oracle}); and rates \textcolor{gray}{\sout{obtained by compressibility}}
\textcolor{colR1}{that follow from the Besov assumption alone, through a bound on the
ideal risk over chunks}
(Theorem~\ref{thm:main})\textcolor{colR1}{, together with a minimax lower bound, for any number
of modulators, that makes these rates optimal when $\pi \ge 2$ (Theorem~S6.1 in
the Supplementary Material)}. \textcolor{gray}{\sout{The second is the reading that compressibility costs
no assumption: the Besov condition on the components already implies that
their wavelet coefficients lie in a weak $\ell_\tau$ ball with
$\tau = (s+1/2)^{-1}$ (Lemma~\mbox{\ref{lem:weak}}, a form with explicit constants
of the classical embedding in \mbox{\citealp[Proposition~9.15]{johnstone2019gaussian}}),
and it is this that separates the rate of the estimator from the rate of linear
approximation when the Besov integrability index is smaller than two.}}
\textcolor{colR1}{The second is the recovery of the structure of the model, which
modulators act on which coefficients, by a threshold on the norms of the
fitted components (Corollary~\ref{cor:threshold}): removing the null blocks
needs no separation condition, and the thresholded estimator keeps the rate of
the fit (Corollary~\ref{cor:thrisk}), so that one estimator serves both for
structure and for prediction.} The third is computational: the
estimator is one convex program \textcolor{gray}{\sout{on a sparse design matrix}} \textcolor{colR1}{solved by
standard group lasso software}, \textcolor{gray}{\sout{and}} the resolution
level and the penalty are chosen jointly \textcolor{colR1}{by cross-validation, and the
threshold by the one standard error rule on the same folds}.
\textcolor{gray}{\sout{The fourth is the recovery of the structure of the model, which
modulators act on which coefficients, by a threshold on the norms of the
fitted components (Corollary~\mbox{\ref{cor:threshold}}): it needs no design
condition beyond the one above and no irrepresentability condition, and
removing the null blocks needs no separation condition at all.}}

The rest of the paper is organized as follows. Section~\ref{sec:model}
introduces the model, its identification, the wavelet \textcolor{colR1}{approximation space} and the estimator.
Section~\ref{sec:theory} states the assumptions and the theory, opening with
the main rate and then giving the three results it rests on\textcolor{colR1}{, then
the rates, including one with no design condition,} \textcolor{colR1}{and
closes with the recovery of the block structure by thresholding} \textcolor{colR1}{and
the risk of the thresholded estimator}.
Section~\ref{sec:computation} describes the computation, the scale of the
penalty and the tuning rules\textcolor{colR1}{, the threshold among them}.
% E5h (k = 5): as Secoes 5 e 6 no roteiro; a Secao 7 (Discussion) e a frase de
% fechamento continuam com a E5b.
\textcolor{colR1}{Section~\ref{sec:simulation} compares the method with penalized
splines and with other competitors in a simulation study, and
Section~\ref{sec:application} applies it to hourly air pollution data from
London.} Proofs are collected in the online Supplementary
Material\textcolor{colR1}{, which also gives the theory of the coordinatewise lasso,
offered by the software as an option}\textcolor{colR1}{, the details and further
results of the simulation study, and a second application, to data from
Beijing}.

Throughout, $\norm{a}_\pi$ is the $\ell_\pi$ norm of a vector,
$\normn{h}^2 = n^{-1}\sum_{i\le n} h_i^2$ is the empirical norm,
$\lmin(\cdot)$ and $\lmaxx(\cdot)$ are the extreme eigenvalues of a symmetric
matrix, $\opnorm{\cdot}$ is the operator norm, and $a_n \asymp b_n$ means that
$a_n/b_n$ is bounded away from zero and infinity. For a support
$S$, $\delta_S$ is the vector that agrees with $\delta$ on $S$ and vanishes
elsewhere.

%%%%%%%%%%%%%%%%%%%%%%%%%%%%%%%%%%%%%%%%%%%%%%%%%%%%%%%%%%%%%%%%%%%%%%%%%%%%%%
\section{Model and Wavelet \texorpdfstring{\textcolor{colR1}{Approximation Space}}{Approximation Space}}
\label{sec:model}

\subsection{The additive functional coefficient model}
\label{sec:model-def}

Let $Y$ be a scalar response, $\bX = (X_1,\dots,X_p)\T$ a vector of linear
covariates and $\bU = (U_1,\dots,U_q)\T$ a vector of modulating covariates,
taking values in $[0,1]^q$ after a monotone transformation of each
coordinate. The model is
\begin{equation}
  \label{eq:model}
  Y \;=\; \sum_{\cv=1}^{p} \fcoef{\cv}(\bU)\,X_{\cv} \;+\; \varepsilon,
  \qquad \E(\varepsilon \mid \bX, \bU) = 0,
\end{equation}
with each functional coefficient additive in the modulators,
\begin{equation}
  \label{eq:additive}
  \fcoef{\cv}(u) \;=\; \lvl{\cv} \;+\; \sum_{\md=1}^{q}\gcomp{\cv}{\md}(u_{\md}),
  \qquad \int_{0}^{1}\gcomp{\cv}{\md}(v)\,dv = 0,
  \qquad \cv = 1,\dots,p .
\end{equation}
The constant $\lvl{\cv}$ is the level of the $\cv$th coefficient and
$\gcomp{\cv}{\md}$ is the component of that coefficient along the $\md$th
modulator. The centering in \eqref{eq:additive} is with respect to Lebesgue
measure, which is the convention of the wavelet literature and the one the
basis of Section~\ref{sec:sieve} imposes at no cost; the parametrization
centered at the law of $\bU$ is obtained from it by the shift
$\lvl{\cv} \mapsto \lvl{\cv} + \sum_{\md}\E\{\gcomp{\cv}{\md}(U_{\md})\}$,
which moves mass between the levels and the components but changes neither
$\fcoef{\cv}$ nor the estimator.

The model contains several familiar ones. With $X_1 \equiv 1$ the first
coefficient is an additive intercept, so that $p = 1$ gives the additive model
and $p > 1$ with $\gcomp{\cv}{\md} \equiv 0$ for $\cv \ge 2$ gives the
partially linear additive model; $q = 1$ gives the classical varying
coefficient model with a single index; and $\gcomp{\cv}{\md}\equiv 0$ for all
pairs gives linear regression. Writing
$f(x,u) = \sum_{\cv} x_{\cv}\fcoef{\cv}(u)$ for the regression function, the
first question is whether the distribution of $(Y,\bX,\bU)$ determines the
parameters. Two degeneracies have to be excluded: if $\bX$ is confined to a
hyperplane given $\bU$, only a combination of the coefficients is determined,
and if one modulator is a function of another, only the sum of the
corresponding components is.

% derivations/01-identificabilidade.md, Proposicao 1 (numeracao global).
\begin{proposition}[Identification]
\label{prop:identif}
Suppose that $\E\norm{\bX}_2^2 < \infty$, that $\E(\bX\bX\T \mid \bU)$ is
nonsingular almost surely, and that $\bU$ has a Lebesgue density on $[0,1]^q$
that is positive almost everywhere. Then, among the parametrizations that
satisfy \eqref{eq:additive}, the levels $\lvl{1},\dots,\lvl{p}$ and the
components $\gcomp{\cv}{\md} \in L_2[0,1]$ are determined by $f$: the levels
uniquely, and each component up to a Lebesgue null set.
\end{proposition}

Dependence between $\bX$ and $\bU$ is allowed, and so is dependence among the
modulators; what the two conditions exclude is a linear combination of the
linear covariates that is a function of the modulators, and a modulator that
is a function of the others. Both exclusions return in quantitative form in
Section~\ref{sec:theory}.

\subsection{The wavelet \texorpdfstring{\textcolor{colR1}{approximation space}}{approximation space}}
\label{sec:sieve}

Let $\phi$ and $\psi$ be the scaling function and the mother wavelet of an
orthonormal multiresolution analysis of $L_2(\R)$, compactly supported, with
$\psi$ bounded and with $N$ vanishing moments, and let $L$ be the length of
the filter, so that the support of \textcolor{gray}{\sout{$\psi$ is contained in $[0,L-1]$}}\textcolor{colR1}{$\phi$ is contained in $[0,L-1]$ and that of $\psi$ in $[1-L/2,\,L/2]$, an interval of length $L-1$} \citep[\textcolor{colR1}{Chapter~7}]{hardle1998wavelets}. Write
\begin{equation}
  \label{eq:periodized}
  \wav{\lev}{\tsl}(u) \;=\; \sum_{\iota \in \mathbb{Z}}
  2^{\lev/2}\psi\{2^{\lev}(u+\iota) - \tsl\},
  \qquad \lev \ge 0, \quad \tsl = 0,\dots,2^{\lev}-1,
\end{equation}
for the periodization of the basis to $[0,1]$. The system
$\{\mathbf{1}\}\cup\{\wav{\lev}{\tsl}\}$ is an orthonormal basis of
$L_2[0,1]$ under Lebesgue measure and every wavelet in it integrates to zero
\citep[Section 9.3]{Daubechies-1992}. We take the coarsest level to be
$\jzero = 0$, so that the only scaling function is the constant
$\phi_{00}\equiv 1$, and we discard it: it is already carried by the level
$\lvl{\cv}$. The \textcolor{colR1}{approximation space (the sieve)} of resolution $J$ is then
\[
  \WJ \;=\; \mathrm{span}\{\wav{\lev}{\tsl} : 0 \le \lev < J,\;
  0 \le \tsl < 2^{\lev}\},
  \qquad \dim \WJ = \NJ = 2^{J}-1,
\]
the orthogonal complement of the constants inside the multiresolution space
$\VJ$, and $\PiJ$ is the orthogonal projection onto it. Every element of $\WJ$
integrates to zero, so that the identification constraint of
\eqref{eq:additive} holds by construction, no numerical restriction is imposed
on the fit, and the \textcolor{colR1}{approximation space} is a plain linear space.

% derivations/01-identificabilidade.md, Lema 1 (numeracao global).
\begin{lemma}[The basis carries the constraint]
\label{lem:basis}
Under the conditions of Proposition~\ref{prop:identif}, for every $J$ the map
\[
  (\bc, \btheta) \;\longmapsto\;
  \sum_{\cv=1}^{p} x_{\cv}\Bigl\{\lvl{\cv} + \sum_{\md=1}^{q}
  \sum_{\lev<J}\sum_{\tsl} \coef{\cv}{\md}{\lev}{\tsl}
  \wav{\lev}{\tsl}(u_{\md})\Bigr\}
\]
is injective on $\R^{p}\times\R^{d}$, $d = pq\NJ$; equivalently, the
population Gram matrix of the design defined in \eqref{eq:design} is positive
definite. If instead the column $X_{\cv}\phi_{00}(U_{\md})$ is retained in
each block, that matrix is singular with nullity exactly $pq$.
\end{lemma}

Replacing $\gcomp{\cv}{\md}$ by $\PiJ\gcomp{\cv}{\md}$ leaves a bias, which
Section~\ref{sec:theory} controls under a Besov condition. Two practical
consequences of periodization deserve mention. A component whose values at the
two ends of the interval disagree is not periodic, and its projection error
then stalls at the order $2^{-J/2}$ in $L_2$ regardless of the interior
smoothness (Proposition~S3.1 in the Supplementary Material). The two remedies,
rescaling the modulators to a proper subinterval and using the basis of
\citet*{cohen1993wavelets} adapted to the interval, are matters of computation
and are discussed in Section~\ref{sec:computation}.

\subsection{The estimator}
\label{sec:estimator}

Let $(Y_i,\bX_i,\bU_i)$, $i = 1,\dots,n$, be independent copies of
$(Y,\bX,\bU)$. The row of the design matrix is the vector of the $p$ linear
covariates followed by the $d = pq\NJ$ products between a linear covariate and
a wavelet of a modulator,
\begin{equation}
  \label{eq:design}
  \bZ = [\,\Amat \;\; \Bmat\,], \qquad
  \Amat_{i\cv} = X_{i\cv}, \qquad
  \Bmat_{i,(\cv\md,\lev\tsl)} = X_{i\cv}\,\wav{\lev}{\tsl}(U_{i\md}),
\end{equation}
with the blocks $(\cv,\md)$ in lexicographic order and, inside a block, the
wavelets ordered by level and then by translation. \textcolor{colR1}{The penalty
acts on chunks of wavelet coefficients inside a block. Fix an integer
$b \ge 2$ and let $\jst$ be the finest level with $2^{\jst} < b$. In each block
$(\cv,\md)$ the levels $\lev \le \jst$ form a single chunk, the coarse chunk,
and each finer level, which has $2^{\lev} \ge b$ coefficients, is cut into
$\lfloor 2^{\lev}/b\rfloor$ chunks of consecutive translations, all with $b$
coefficients except the last, which also takes the remainder. Write $\Gcal$ for
the resulting partition of the $d$ penalized coordinates, $\btheta_{G}$ and
$|G|$ for the subvector and the size of a chunk $G$, and, for positive weights
$w = (w_{G})_{G\in\Gcal}$, $\nGw{\btheta} = \sum_{G\in\Gcal}w_{G}\norm{\btheta_{G}}_{2}$.}
Writing
$\bY = (Y_1,\dots,Y_n)\T$, the estimator is
% derivations/08-blocos.tex, eq. (est-blocos) e a forma balanceada (D44).
{\color{gray}\removedk
\[
  (\chat,\thetahat) \;\in\; \mathop{\rm arg\,min}_{(\bc,\btheta)\in
  \R^{p}\times\R^{d}}
  \Bigl\{\normn{\bY - \Amat\bc - \Bmat\btheta}^{2}
  \;+\; 2\pen\norm{\btheta}_{1}\Bigr\} .
\]
}
{\color{colR1}
\begin{equation}
  \label{eq:estimator}
  (\chat,\thetahat) \;\in\; \mathop{\rm arg\,min}_{(\bc,\btheta)\in
  \R^{p}\times\R^{d}}
  \Bigl\{\normn{\bY - \Amat\bc - \Bmat\btheta}^{2}
  \;+\; 2\pen\nGw{\btheta}\Bigr\} ,
\end{equation}
with $b = \bn = \lceil\log n\rceil$ and $w_{G} = |G|^{1/2}$. Every chunk of a
fine level has between $b$ and $2b-1$ coefficients, the coarse chunk has
$2^{\jst+1}-1$, which lies between $b-1$ and $2b-3$ (when $J \le \jst+1$
each block is a single chunk), there are at most
$pq(1+2^{J}/b)$ chunks, and the ratio $\varrho = \max_{G}w_{G}/\min_{G}w_{G}$
of the weights satisfies $\varrho^{2} \le (2b-1)/(b-1) \le 3$
(Proposition~S5.1 in the Supplementary Material). The theory of
Section~\ref{sec:theory} holds for any weights whose ratio stays bounded, and
the bound on $\varrho$ is what makes it cover the weights $|G|^{1/2}$ that
group lasso software uses by default.}
The penalty does not reach the levels $\bc$: they are $p$ parameters of a
parametric part, and shrinking them would bias every fitted coefficient by a
constant. The fitted components and coefficients are recovered from
$\thetahat$ by
\[
  \widehat{g}_{\cv\md}(v) = \sum_{\lev<J}\sum_{\tsl}
  \widehat{\theta}_{\cv\md,\lev\tsl}\,\wav{\lev}{\tsl}(v),
  \qquad
  \widehat{\fcoef{\cv}}(u) = \widehat{c}_{\cv}
  + \sum_{\md=1}^{q}\widehat{g}_{\cv\md}(u_{\md}).
\]
% E5f (L10; D54): a citacao de Yuan & Lin dentro do grupo azul de k = 4.
Problem \eqref{eq:estimator} is a \textcolor{gray}{\sout{lasso}} \textcolor{gray}{\sout{\mbox{\citep{tibshirani1996regression}}}}\textcolor{colR1}{group lasso \citep{Yuan-Lin-2006}} with some coordinates unpenalized, so
it is convex, and it is strictly convex, hence uniquely solved, whenever the
Gram matrix of $\bZ$ is nonsingular. \textcolor{gray}{\sout{Two structural variants are available at
no conceptual cost and are discussed in Section~\mbox{\ref{sec:computation}}: a
sparse group lasso \mbox{\citep{Simon-Friedman-Hastie-Tibshirani-2013}} whose groups
are the blocks $(\cv,\md)$, which removes an entire component when a modulator
does not act on a coefficient, and the block lasso of
\mbox{\citet{Klopp-Pensky-2015}}.}} \textcolor{colR1}{It differs from the block lasso of
\citet{Klopp-Pensky-2015} in three respects: their chunks run across resolution
levels, their weights are equal, and their levels are penalized with everything
else. Two other penalties are available and are discussed in
Section~\ref{sec:computation}: the lasso \citep{tibshirani1996regression}, with
$\norm{\btheta}_{1}$ in place of $\nGw{\btheta}$, whose theory is given in
Section~S8 of the Supplementary Material, and a sparse group lasso
\citep{Simon-Friedman-Hastie-Tibshirani-2013} whose groups are the blocks
$(\cv,\md)$.} The theory below is for \eqref{eq:estimator}.

%%%%%%%%%%%%%%%%%%%%%%%%%%%%%%%%%%%%%%%%%%%%%%%%%%%%%%%%%%%%%%%%%%%%%%%%%%%%%%
\section{Theory}
\label{sec:theory}

\subsection{Assumptions}
\label{sec:assumptions}

\begin{assumption}[Sampling and error]
\label{ass:sample}
The observations $(Y_i,\bX_i,\bU_i)$, $i \le n$, are independent and
identically distributed and follow \eqref{eq:model}, and the error is
conditionally sub-Gaussian: there is $\sigma < \infty$ with
$\E\{\exp(t\varepsilon_i) \mid \bX_i,\bU_i\} \le \exp(\sigma^2t^2/2)$ almost
surely, for all real $t$.
\end{assumption}

\begin{assumption}[Linear covariates]
\label{ass:X}
$\max_{\cv}|X_{\cv}| \le B_X$ almost surely, and there are constants
$0 < \kappa_1 \le \kappa_2 < \infty$ with
$\kappa_1 \le \lmin\{\E(\bX\bX\T\mid\bU)\} \le
\lmaxx\{\E(\bX\bX\T\mid\bU)\} \le \kappa_2$ almost surely.
\end{assumption}

\begin{assumption}[Modulating covariates]
\label{ass:U}
$\bU$ takes values in $[0,1]^q$ and has a joint Lebesgue density $p_{\bU}$
with $0 < c_U \le p_{\bU}(u) \le C_U < \infty$ for every $u \in [0,1]^q$.
\end{assumption}

\begin{assumption}[Basis]
\label{ass:basis}
The wavelets are the periodization \eqref{eq:periodized} of an orthonormal,
compactly supported family with $\psi$ bounded and $N$ vanishing moments, and
$\jzero = 0$.
\end{assumption}

\begin{assumption}[Besov regularity]
\label{ass:besov}
There are $s > 0$, $1 \le \pi, r \le \infty$ and $C_g < \infty$ such that
every component satisfies $\int_0^1 \gcomp{\cv}{\md} = 0$ and, writing
$\coefs{\cv}{\md}{\lev}{\tsl} = \inner{\gcomp{\cv}{\md}}{\wav{\lev}{\tsl}}$
and $\theta^{*}_{\cv\md,\lev\cdot}$ for the vector of the coefficients of
level $\lev$,
\[
  \norm{\theta^{*}_{\cv\md,\lev\cdot}}_{\pi} \;\le\;
  C_g\,2^{-\lev(s + 1/2 - 1/\pi)} \qquad \text{for every } \lev \ge 0 .
\]
The effective regularity $\seff = s - (1/\pi - 1/2)_{+}$ is positive.
\end{assumption}

Assumptions~\ref{ass:X} and \ref{ass:U} are the quantitative forms of the two
conditions of Proposition~\ref{prop:identif}. In the second it is the joint
density that is bounded below, not the marginals: if $U_2 = U_1$ almost surely
both marginals are uniform and the design is singular.
Assumption~\ref{ass:besov} is taken as a condition on the sequence of
coefficients rather than as membership in a function space, as in
\citet{Klopp-Pensky-2015}; when $\psi$ has more than $s$ vanishing moments and
enough Hölder regularity the two are equivalent, on the circle
\citep[Section 3.6]{Cohen-2003}, and the implication that the proofs use needs
only the vanishing moments. The index $\seff$ is the regularity that \textcolor{colR1}{linear
approximation} sees and $s$ is the one that nonlinear approximation sees; they differ
exactly when $\pi < 2$, which is the regime of inhomogeneous smoothness.

\subsection{The main result}
\label{sec:main}

Write $\widehat{f}_i = \Amat_i\chat + \Bmat_i\thetahat$ for the fit, \textcolor{gray}{\sout{$d_n$ for
the number of penalized columns at resolution $J_n$}}\textcolor{colR1}{$\Gcal_{n}$ for the
partition of Section~\ref{sec:estimator} at resolution $J_n$ with $b = \bn$},
and fix a confidence
level $\alpha \in (0,1)$. The penalty is the one that the oracle inequality of
Section~\ref{sec:oracle} calls for, namely\textcolor{colR1}{, with $\gamma = \kappa_1c_U$
and $\Lambda = \kappa_2C_U + \gamma/2$ the constants of
Proposition~\ref{prop:gram} below,}
{\color{gray}\removedk
\[
  \penn \;=\; 2\sigma B_X\sqrt{2C_U}\,
  \sqrt{\frac{2\log(2d_n/\alpha)}{n}} \;\asymp\;
  \sigma B_X \sqrt{C_U}\,\sqrt{\frac{J_n}{n}} .
\]
}
% derivations/08-blocos.tex, Lema 14(ii)-(iii) e a Hipotese B.
{\color{colR1}
\begin{equation}
  \label{eq:lambda}
  \lamG \;=\; 2\max_{G\in\Gcal_{n}}\frac{\lamzbar{G}}{w_{G}},
  \quad
  \lamzbar{G} \;=\; \frac{\sigma}{\sqrt{n}}\Bigl\{B_X\sqrt{2C_U|G|}
  + \sqrt{2\Lambda\log(|\Gcal_{n}|/\alpha)}\Bigr\} .
\end{equation}
For every chunk, $\lamG w_{G} \le \varrho\lamGone$, where
$\lamGone = 2\max_{G}\lamzbar{G}$ is of order $\sigma(\log n/n)^{1/2}$, the
order of the penalty of the lasso.}

% derivations/05-taxas.tex, Corolario 5 (numeracao global); D16. Substituido
% em k = 4 pelo Corolario 11 de derivations/08-blocos.tex (D47(a)).
{\color{gray}\removedk
\begin{theoremold}[Rate of nonlinear approximation]
Let Assumptions~\ref{ass:sample} to \ref{ass:besov} hold with $p$, $q$ and all
the constants fixed as $n$ grows, and set $\tau = (s+1/2)^{-1}$. Suppose
$\seff > s/(2s+1)$ and take $J_n = \lceil c \log_2 n\rceil$ with
\[
  \max\Bigl\{\frac{\tau}{2},\; \frac{2-\tau}{4\seff}\Bigr\}
  \;\le\; c \;<\; 1 ,
\]
a window that is nonempty precisely under that condition. Then, with $\penn$
as in the display above,
\[
  \normn{\widehat{f}-f}^{2}
  \;=\; O_p\Bigl\{\Bigl(\frac{\log n}{n}\Bigr)^{2s/(2s+1)}\Bigr\},
  \qquad
  \sum_{\cv,\md}\bigl\lVert \widehat{g}_{\cv\md}-\gcomp{\cv}{\md}
  \bigr\rVert_{L_2[0,1]}^{2}
  \;=\; O_p\Bigl\{\Bigl(\frac{\log n}{n}\Bigr)^{2s/(2s+1)}\Bigr\} .
\]
\end{theoremold}
}
% derivations/08-blocos.tex, Corolario 11 (numeracao global); D47(a).
{\color{colR1}
\begin{theorem}[Rate of nonlinear approximation]
\label{thm:main}
Let Assumptions~\ref{ass:sample} to \ref{ass:besov} hold with $p$, $q$ and all
the constants fixed as $n$ grows, and let the weights have ratio
$\varrho \le \bar\varrho$ for a constant $\bar\varrho$ fixed as $n$ grows
($\bar\varrho^{2} = 3$ for the weights of Section~\ref{sec:estimator}). Suppose
$\seff > s/(2s+1)$ and take $J_n = \lceil c \log_2 n\rceil$ with
\begin{equation}
  \label{eq:window}
  \frac{s}{(2s+1)\seff} \;\le\; c \;<\; 1 ,
\end{equation}
a window that is nonempty precisely under that condition. Then, with
$\pen = \lamG$ as in \eqref{eq:lambda} and
$r_n = n^{-2s/(2s+1)}(\log n)^{(2/\pi-1)_{+}/(2s+1)}$,
\[
  \normn{\widehat{f}-f}^{2} \;=\; O_p(r_n),
  \qquad
  \sum_{\cv,\md}\bigl\lVert \widehat{g}_{\cv\md}-\gcomp{\cv}{\md}
  \bigr\rVert_{L_2[0,1]}^{2} \;=\; O_p(r_n) .
\]
In particular, when $\pi \ge 2$ the rate is $n^{-2s/(2s+1)}$, with no
logarithmic factor.
\end{theorem}
}

Three readings of Theorem~\ref{thm:main} matter. First, the exponent is the
one of nonlinear wavelet approximation, $2s/(2s+1)$, and not the one of
\textcolor{colR1}{linear approximation}, $2\seff/(2\seff+1)$ of Theorem~\ref{thm:sieve} below; the two
agree when $\pi \ge 2$ and separate when $\pi < 2$, which is where the
\textcolor{gray}{\sout{$\ell_1$}} penalty pays for itself. Second, the rate is attained under
Assumption~\ref{ass:besov} alone: \textcolor{gray}{\sout{no compressibility is assumed on top of it,
because the Besov condition already implies it, with
$\tau = (s+1/2)^{-1}$, by Lemma~\mbox{\ref{lem:weak}}.}} \textcolor{colR1}{no condition is imposed
on the coefficients on top of it, because the Besov condition already bounds
the ideal risk over chunks that governs the estimation error
(Lemma~\ref{lem:weak}).} \textcolor{gray}{\sout{Third, the rate
$(\log n/n)^{2s/(2s+1)}$ is the minimax reference of the Gaussian sequence
model over Besov bodies with per coordinate noise of order $\penn$
(\mbox{\citealp[Theorems~4 and~5]{Donoho-Johnstone-1998}};
\mbox{\citealp[Theorem~9.3]{johnstone2019gaussian}}); the logarithmic factor is the price of the
single, nonadaptive penalty level \mbox{\eqref{eq:lambda}}, which takes a union bound
over the $d_n$ columns. No lower bound is proved here for $q \ge 2$: the
comparison is with the reference of the sequence model, and the minimax lower
bound of \mbox{\citet{Klopp-Pensky-2015}} covers the case of a single modulator with
covariates independent of it.}}
% E5g (D57, pergunta 51(b)): a terceira leitura passa a citar a cota inferior
% do Theorem S6.1 do supp (Teorema 5 de derivations/09-cota-inferior.tex).
\textcolor{colR1}{Third, \textcolor{gray}{\sout{$n^{-2s/(2s+1)}$ is the minimax rate over Besov bodies
in the Gaussian sequence model with noise level of order $n^{-1/2}$
(\mbox{\citeauthor{Donoho-Johnstone-1998}}, \mbox{\citeyear{Donoho-Johnstone-1998}}, Theorems~4 and~5), and with a single modulator,
independent of the covariates, and Gaussian errors, the minimax lower bound of
\mbox{\citeauthor{Klopp-Pensky-2015}} (\mbox{\citeyear{Klopp-Pensky-2015}}, Theorem~1) holds over the class of
\mbox{Assumption~\ref{ass:besov}} (Remark~S6.1 in the Supplementary Material), so that
in that case \mbox{Theorem~\ref{thm:main}} is minimax optimal when $\pi \ge 2$.}} the rate
cannot be improved. With Gaussian errors and any law of $(\bX,\bU)$ that
satisfies Assumptions~\ref{ass:X} and~\ref{ass:U}, the minimax risk over the
class of Assumption~\ref{ass:besov}, in the prediction norm and for the
components, is at least a constant multiple of $pq\,n^{-2s/(2s+1)}$, for every
$\pi$ (Theorem~S6.1 in the Supplementary Material). This is the regression
counterpart of the lower bound of \citet[Theorem~5]{Donoho-Johnstone-1998} for
Besov bodies in the Gaussian sequence model, and with a fixed number of
covariates it extends that of \citet[Theorem~1]{Klopp-Pensky-2015}, who have a
single modulator independent of the covariates. Theorem~\ref{thm:main} is
therefore minimax optimal when $\pi \ge 2$, with any number of modulators and
with covariates that depend on them. When
$\pi < 2$ the remaining factor $(\log n)^{(2/\pi-1)/(2s+1)}$ is the one in the
rate of \citet[eq.~3.19]{Klopp-Pensky-2015} with a fixed number of covariates,
and in that of block thresholding in the sequence model
\citep[Theorem~4]{Cai-1999}; \textcolor{gray}{\sout{whether it is necessary for this estimator is not
known, and no lower bound is proved here for $q \ge 2$.}} the lower bound carries
no such factor, and whether the factor is necessary for this estimator is not
known. Compared with the
coordinatewise lasso, whose rate under the same conditions is
$(\log n/n)^{2s/(2s+1)}$ (Theorem~S8.3 in the Supplementary Material), the
chunks remove a factor $(\log n)^{2\seff/(2s+1)}$.
% D55 (pergunta 49(a)): a frase de D49(c) passa a valer para todo pi, pela S6.2.
\textcolor{gray}{\sout{When $s < 1/2$ and
$\pi \ge 2$, the rate $n^{-2s/(2s+1)}$ of \mbox{Theorem~\ref{thm:main}} is attained
with no design condition at all, at another resolution
\mbox{(Proposition~\ref{prop:slow}(ii) below)}.}} When $s < 1/2$, the prediction
rate of Theorem~\ref{thm:main} holds with no design condition at all, under
its other assumptions and in the same window for $J_{n}$ (Proposition~S6.2
in the Supplementary Material).}

The proof combines three ingredients, given below: the approximation error of
\textcolor{colR1}{$\WJ$}, the design condition, and a nonasymptotic oracle inequality.
Section~\ref{sec:rates} assembles them, first at \textcolor{gray}{\sout{a dense reading of the
estimation term}}\textcolor{colR1}{the oracle}, which gives the \textcolor{colR1}{rate of linear approximation}, and then at \textcolor{gray}{\sout{a sparse comparator}}\textcolor{colR1}{a
comparator that keeps only the chunks of large energy},
which gives Theorem~\ref{thm:main}.

\subsection{Approximation and the design condition}
\label{sec:approx}

% derivations/02-aproximacao-besov.tex, Lema 2 e Corolario 1 (numeracao global).
\begin{lemma}[\textcolor{colR1}{Approximation bias}]
\label{lem:bias}
Under Assumptions~\ref{ass:basis} and \ref{ass:besov}, for every pair
$(\cv,\md)$ and every $J \ge 0$,
\[
  \bigl\lVert \gcomp{\cv}{\md} - \PiJ\gcomp{\cv}{\md}
  \bigr\rVert_{L_2[0,1]}^{2}
  \;=\; \sum_{\lev \ge J}\norm{\theta^{*}_{\cv\md,\lev\cdot}}_{2}^{2}
  \;\le\; \frac{C_g^{2}}{1-2^{-2\seff}}\,2^{-2J\seff} .
\]
If in addition Assumptions~\ref{ass:X} and \ref{ass:U} hold, then, with
$f_J(x,u) = \sum_{\cv}x_{\cv}\{\lvl{\cv} +
\sum_{\md}(\PiJ\gcomp{\cv}{\md})(u_{\md})\}$,
\begin{equation}
  \label{eq:bias}
  \bigl\lVert f - f_J\bigr\rVert_{L_2(P)}^{2} \;\le\;
  \frac{(pq)^{2}B_X^{2}C_UC_g^{2}}{1-2^{-2\seff}}\,2^{-2J\seff}
  \;=:\; \mathcal{B}_{J}^{2} .
\end{equation}
\end{lemma}

The exponent is $\seff$ and not $s$, and this is where the loss
$(1/\pi-1/2)_{+}$ enters: inside a level, an $\ell_\pi$ bound with $\pi < 2$
controls the $\ell_2$ norm without gain. The bias uses only the bound on the
linear covariates and the marginal density of each modulator, so dependence
between $\bX$ and $\bU$ does not affect it. What it does affect is the passage
from the prediction norm to the norm of each component, which is the design
condition.

Let $\bPsi(u) \in \R^{K}$, $K = 1 + q\NJ$, collect the constant and the
wavelets of the $q$ modulators, so that the row of $\bZ$ in
\eqref{eq:design} is a permutation of the Kronecker product
$\bX \otimes \bPsi(\bU)$. Write $\Gram = \E(\bZ\bZ\T)$ and
$\Gramhat = \bZ\T\bZ/n$ for the population and empirical Gram matrices and
$\Gram_{\bPsi} = \E\{\bPsi(\bU)\bPsi(\bU)\T\}$ for the part in the modulators
alone.

% derivations/03-desenho-produtos.tex, Proposicao 3 (numeracao global).
\begin{proposition}[Gram matrix of the product design]
\label{prop:gram}
Let Assumptions~\ref{ass:X} to \ref{ass:basis} hold and write $D = p + d$ and
$\gamma = \kappa_1 c_U$.
\begin{enumerate}
\item[(i)] $c_U \le \lmin(\Gram_{\bPsi}) \le \lmaxx(\Gram_{\bPsi}) \le C_U$
for every $J$, and consequently
$\gamma \le \lmin(\Gram) \le \lmaxx(\Gram) \le \kappa_2 C_U$. If $\bX$ and
$\bU$ are independent, $\Gram$ is, up to the permutation of columns, the
Kronecker product $\E(\bX\bX\T)\otimes\Gram_{\bPsi}$, and the two bounds are
attained as products of eigenvalues.
\item[(ii)] With $R_J = pB_X^{2}(1 + qC_\psi 2^{J})$, where $C_\psi$ depends
only on the wavelet family, $\norm{\bZ_i}_2^{2} \le R_J$ almost surely and
\[
  \Prob\Bigl(\lmin(\Gramhat) \ge \frac{\gamma}{2}\Bigr)
  \;\ge\; 1 - 2D\exp\Bigl\{-\frac{n\gamma^{2}}
  {8R_J(\kappa_2C_U + \gamma/3)}\Bigr\},
\]
which tends to one whenever $pq2^{J}\log(pq2^{J})/n \to 0$.
\end{enumerate}
\end{proposition}

Part (i) is stated for the general case: independence between the covariates
and the modulators buys the exact factorization but not a better constant. A
lower bound on the smallest eigenvalue dominates every compatibility and
restricted eigenvalue constant\textcolor{colR1}{, including the versions for groups
of \citet{Lounici-Pontil-vandeGeer-Tsybakov-2011}}, on every support and with no cone restriction,
since $\delta\T A \delta \ge \lmin(A)\norm{\delta}_2^2 \ge
\lmin(A)\norm{\delta_S}_1^2/|S|$\textcolor{colR1}{, and likewise
$\delta\T A\delta \ge \lmin(A)\{\sum_{G\in\Hcal}w_{G}\norm{\delta_{G}}_{2}\}^{2}/\Wcal(\Hcal)$
for a set $\Hcal$ of chunks, with $\Wcal(\Hcal) = \sum_{G\in\Hcal}w_{G}^{2}$},
and it is in this form that the next
subsection consumes it. In the resolution regime of Section~\ref{sec:rates}
the requirement of part (ii) is \textcolor{gray}{\sout{of the same order as the rate itself}}\textcolor{colR1}{the
rate times $\log n$, which tends to zero}, so it
costs nothing there; if $pq$ were larger than $n^{1/2}$, only the
compatibility constant could be transferred, through the entrywise deviation
\citep[Corollary 6.8]{buhlmann2011statistics}, at the price of a factor
$s_0^{2}$ and of the cone.

\subsection{A nonasymptotic oracle inequality}
\label{sec:oracle}

Let $\betastar = (\bc,\thetastar)$ be the \textcolor{colR1}{oracle of the approximation space}, that is the vector
whose components are the levels and the coefficients
$\coefs{\cv}{\md}{\lev}{\tsl}$ of Assumption~\ref{ass:besov} for $\lev < J$,
so that $f_J$ is its fit, and let
{\color{gray}\removedk
\[
  \bb = f - f_J \in \R^{n}, \qquad \Szero = \mathrm{supp}(\thetastar),
  \qquad s_0 = |\Szero|, \qquad v = \thetahat - \thetastar .
\]
}
{\color{colR1}
\[
  \bb = f - f_J \in \R^{n},
  \qquad
  \bbar = f - \bZ\bar\bbeta, \qquad \vbar = \thetahat - \thbar ,
  \qquad \Gbar = \{G : \thbar_{G} \ne 0\},
\]
for a comparator $\bar\bbeta = (\bar\bc,\thbar) \in \R^{p}\times\R^{d}$, any
vector of levels and coefficients, $\Gbar$ being its active chunks, and
$\bar\Wcal = \Wcal(\Gbar) = \sum_{G\in\Gbar}w_{G}^{2}$.}
Let $\PA$ be the orthogonal projection onto the column space of $\Amat$, which
exists whenever $\Amat$ has rank $p$, let $\MA = I_n - \PA$, and let
$\Btil = \MA\Bmat$ and $\gtil = \lmin(\Btil\T\Btil/n)$ be the residualized
design and its smallest eigenvalue\textcolor{colR1}{, and
$\Psitil_{G} = \Btil_{G}\T\Btil_{G}/n$ the residualized Gram matrix of the
columns of a chunk}. Finally let
$\smax^{2} = \max_a (\Gramhat)_{aa}$ be the largest empirical norm of a
penalized column.

% derivations/04-oraculo.tex, Teorema 1 (numeracao global). Substituido em
% k = 4 pelo Teorema 3 de derivations/08-blocos.tex; o do lasso e o
% Theorem S8.1 do supp.
{\color{gray}\removedk
\begin{theoremold}[Oracle inequality]
Let Assumption~\ref{ass:sample} hold, suppose $\Amat$ has rank $p$, and take
$\pen \ge 2\pen_0$ with
$\pen_0 = \sigma\smax\{2\log(2d/\alpha)/n\}^{1/2}$. On the event
$\mathcal{T} = \{\norm{\Btil\T\bveps/n}_{\infty} \le \pen/2\}$, which has
probability at least $1 - \alpha$:
\begin{enumerate}
\item[(i)] with no condition on the design,
$\normn{\Btil v}^{2} \le 6\pen\norm{\thetastar}_{1} + 4\normn{\bb}^{2}$;
\item[(ii)] if $\gtil > 0$,
\[
  \normn{\Btil v}^{2} \le \frac{64\pen^{2}s_0}{\gtil} + 16\normn{\bb}^{2},
  \qquad
  \norm{v}_{2}^{2} \le \frac{64\pen^{2}s_0}{\gtil^{2}}
  + \frac{16\normn{\bb}^{2}}{\gtil} ;
\]
\item[(iii)] in either regime,
$\normn{\widehat{f}-f} \le \normn{\Btil v} + 2\normn{\bb}
+ \normn{\PA\bveps}$, and
$\E(\normn{\PA\bveps}^{2}\mid\bX,\bU) \le \sigma^{2}p/n$.
\end{enumerate}
\end{theoremold}
}
% derivations/08-blocos.tex, Teorema 3 e Lema 14(i) (numeracao global).
{\color{colR1}
\begin{theorem}[Oracle inequality]
\label{thm:oracle}
Let Assumption~\ref{ass:sample} hold, suppose $\Amat$ has rank $p$, let the
weights not depend on the responses, and, for $\alpha \in (0,1)$, let
\[
  \lamz{G} \;=\; \frac{\sigma}{\sqrt{n}}\Bigl[
  \{\mathrm{tr}(\Psitil_{G})\}^{1/2}
  + \bigl\{2\opnorm{\Psitil_{G}}\log(|\Gcal|/\alpha)\bigr\}^{1/2}\Bigr]
\]
and $\pen_{w} = 2\max_{G\in\Gcal}\lamz{G}/w_{G}$, and take $\pen \ge \pen_{w}$. On the event
$\Tw = \{\max_{G}\norm{(\Btil\T\bveps/n)_{G}}_{2}/w_{G} \le \pen/2\}$, which
has probability at least $1 - \alpha$, the following hold simultaneously for
every comparator $\bar\bbeta$:
\begin{enumerate}
\item[(i)] with no condition on the design,
$\tfrac12\normn{\Btil\vbar}^{2} + \pen\nGw{\thetahat} \le
3\pen\nGw{\thbar} + 2\normn{\bbar}^{2}$;
\item[(ii)] if $\gtil > 0$,
\[
  \normn{\Btil\vbar}^{2} \le \frac{64\pen^{2}\bar\Wcal}{\gtil} + 16\normn{\bbar}^{2},
  \qquad
  \norm{\vbar}_{2}^{2} \le \frac{64\pen^{2}\bar\Wcal}{\gtil^{2}}
  + \frac{16\normn{\bbar}^{2}}{\gtil} ;
\]
\item[(iii)] in either regime,
$\normn{\widehat{f}-f} \le \normn{\Btil\vbar} + 2\normn{\bbar}
+ \normn{\PA\bveps}$, and
$\E(\normn{\PA\bveps}^{2}\mid\bX,\bU) \le \sigma^{2}p/n$.
\end{enumerate}
\end{theorem}
}

\textcolor{gray}{\sout{Three}}\textcolor{colR1}{Four} features separate Theorem~\ref{thm:oracle} from the standard treatment
\citep[Chapter 6]{buhlmann2011statistics} \textcolor{colR1}{and from that of the group
lasso by \citet{Lounici-Pontil-vandeGeer-Tsybakov-2011}}. The unpenalized levels are not
placed inside the support, which would require a compatibility condition on a
cone; they are profiled out, and the Schur complement transfers the full
smallest eigenvalue of $\Gramhat$ to the residualized block,
$\gtil \ge \lmin(\Gramhat)$, so that Proposition~\ref{prop:gram}(ii) is used
as it stands and the constant does not depend on \textcolor{gray}{\sout{$\Szero$}}\textcolor{colR1}{$\Gbar$}. The model is
misspecified inside the \textcolor{colR1}{approximation space}, and the deterministic residual \textcolor{gray}{\sout{$\bb$}}\textcolor{colR1}{$\bbar$} is handled
by Young's inequality rather than by duality, which keeps the statement
nonasymptotic and valid also when the bias dominates\textcolor{colR1}{, and the
statement holds for every comparator at once, which is what
Section~\ref{sec:rates} uses}. \textcolor{gray}{\sout{The penalty is
calibrated by the empirical column norms and not by $\max_{i,a}|Z_{ia}|$: the
former is bounded in probability by $2B_X^2C_U$, while the latter is of order
$2^{J/2}$ and would inflate $\pen$ by that factor and destroy the rates of
Section~\mbox{\ref{sec:rates}}.}} \textcolor{colR1}{The penalty is calibrated chunk by
chunk, by the tail inequality for quadratic forms of
\citet[Theorem~2.1]{Hsu-Kakade-Zhang-2012} with a union bound over the chunks
(Lemma~S5.2 in the Supplementary Material); the event $\Tw$ does not depend on
the weights, which only choose the chunk that dictates $\pen_{w}$, and the
price of unequal weights is the factor $\varrho$ in
$\pen_{w}w_{G} \le 2\varrho\max_{G'}\lamz{G'}$, paid in the estimation term.
Since $\mathrm{tr}(\Psitil_{G}) \le |G|\smax^{2}$ and
$\opnorm{\Psitil_{G}} \le \lmaxx(\Gramhat)$, the deterministic level
\eqref{eq:lambda} dominates $\pen_{w}$ whenever $\smax^{2} \le 2B_X^2C_U$ and
$\lmaxx(\Gramhat) \le \Lambda$, two events of probability tending to one in
the regime of Section~\ref{sec:rates}; it is the empirical column norms that
enter, and not $\max_{i,a}|Z_{ia}|$, which is of order $2^{J/2}$ and would
destroy the rates. The software standardizes the columns and orthonormalizes
each chunk before penalizing it; the theorem covers that estimator as well,
with $\Psitil_{G}$ replaced by $Q_{G}^{-1/2}\Psitil_{G}Q_{G}^{-1/2}$, $Q_{G}$
being the centered Gram matrix of the chunk, which makes $\lamz{G}$ pivotal,
and with $\bar\Wcal$ multiplied by $\max_{G}\lmaxx(Q_{G})$ (Remark~S5.2 in the
Supplementary Material).} The term $\sigma^{2}p/n$ in part (iii) is the price
of leaving the $p$ levels out of the penalty; no choice of $\pen$ removes it,
it is of smaller order than the estimation term whenever
\textcolor{gray}{\sout{$p \lesssim s_0\log(pq\NJ)$}}\textcolor{colR1}{$p \lesssim |\Gbar|\log n$}, and it is absent from
\citet{Klopp-Pensky-2015}, where the levels are penalized along with
everything else.

\subsection{Rates}
\label{sec:rates}

Under Assumption~\ref{ass:besov} the \textcolor{colR1}{oracle} is generically dense: the
coefficients decay but do not vanish, so that \textcolor{gray}{\sout{$s_0 = d$}}\textcolor{colR1}{every chunk is
active} and part (ii) of
Theorem~\ref{thm:oracle}, read at the oracle, delivers the rate of
\textcolor{colR1}{linear approximation}\textcolor{colR1}{, now with no logarithmic factor}.

% derivations/05-taxas.tex, Teorema 2 (numeracao global). Substituido em k = 4
% pelo Teorema 4 de derivations/08-blocos.tex; o do lasso e o Theorem S8.2.
{\color{gray}\removedk
\begin{theoremold}[Rate of linear approximation]
Let Assumptions~\ref{ass:sample} to \ref{ass:besov} hold with $p$, $q$ and the
constants fixed, and take
\[
  J_n = \min\bigl\{J \in \mathbb{N} :\;
  2^{J} \ge (n/\log n)^{1/(2\seff+1)}\bigr\},
\]
so that $J_n = \log_2 n/(2\seff+1) + O(\log\log n)$. Then the requirement of
Proposition~\ref{prop:gram}(ii) holds along $J_n$, being itself of the exact
order $(\log n/n)^{2\seff/(2\seff+1)}$, and, with $\penn$ as in
the first display of Section~\ref{sec:main},
\[
  \normn{\widehat{f}-f}^{2}
  \;=\; O_p\Bigl\{\Bigl(\frac{\log n}{n}\Bigr)^{2\seff/(2\seff+1)}\Bigr\} .
\]
At this resolution the estimation term and the bias term are of the same
order, and the term $\sigma^{2}p/n$ of Theorem~\ref{thm:oracle}(iii) is of
strictly smaller order.
\end{theoremold}
}
% derivations/08-blocos.tex, Teorema 4 (numeracao global).
{\color{colR1}
\begin{theorem}[Rate of linear approximation]
\label{thm:sieve}
Let Assumptions~\ref{ass:sample} to \ref{ass:besov} hold with $p$, $q$ and the
constants fixed, let the weights be as in Theorem~\ref{thm:main}, and take
\begin{equation}
  \label{eq:Jn}
  J_n = \min\bigl\{J \in \mathbb{N} :\;
  2^{J} \ge n^{1/(2\seff+1)}\bigr\}.
\end{equation}
Then the requirement of Proposition~\ref{prop:gram}(ii) holds along $J_n$,
being of order $n^{-2\seff/(2\seff+1)}\log n$, and, with $\pen = \lamG$ as in
\eqref{eq:lambda},
\[
  \normn{\widehat{f}-f}^{2}
  \;=\; O_p\bigl(n^{-2\seff/(2\seff+1)}\bigr) .
\]
At this resolution the estimation term, of order $pq(\bn + 2^{J_n})/n$, and
the bias term are of the same order, and the term $\sigma^{2}p/n$ of
Theorem~\ref{thm:oracle}(iii) is of strictly smaller order.
\end{theorem}
}

% derivations/05-taxas.tex, Corolario 4 (numeracao global). Substituido em
% k = 4 pelo Corolario 10 de derivations/08-blocos.tex.
{\color{gray}\removedk
\begin{corollaryold}[Components and coefficients]
In the setting of Theorem~\ref{thm:sieve}, and with $r_n$ denoting the rate of
that theorem,
\[
  \sum_{\cv,\md}\bigl\lVert\widehat{g}_{\cv\md}-\gcomp{\cv}{\md}
  \bigr\rVert_{L_2[0,1]}^{2} = O_p(r_n),
  \qquad
  \bigl\lVert\widehat{\fcoef{\cv}}-\fcoef{\cv}
  \bigr\rVert_{L_2([0,1]^{q})} = O_p(r_n^{1/2}),
  \quad \cv = 1,\dots,p,
\]
and the same statements hold with the rate of Theorem~\ref{thm:main} in place
of $r_n$ under the conditions of that theorem.
\end{corollaryold}
}
% derivations/08-blocos.tex, Corolario 10 (numeracao global).
{\color{colR1}
\begin{corollary}[Components and coefficients]
\label{cor:components}
In the setting of Theorem~\ref{thm:sieve}, with
$\rhoG = n^{-\seff/(2\seff+1)}$,
\[
  \sum_{\cv,\md}\bigl\lVert\widehat{g}_{\cv\md}-\gcomp{\cv}{\md}
  \bigr\rVert_{L_2[0,1]}^{2} = O_p\bigl\{(\rhoG)^{2}\bigr\},
  \qquad
  \bigl\lVert\widehat{\fcoef{\cv}}-\fcoef{\cv}
  \bigr\rVert_{L_2([0,1]^{q})} = O_p(\rhoG)
\]
for $\cv = 1,\dots,p$, and $\norm{\chat - \bc}_{2} = O_p(\rhoG)$.
\end{corollary}
}

The passage from the prediction norm to the norm of a single component is
where the design condition is used a second time, and the smallest eigenvalue
appears squared in the constant. \textcolor{colR1}{In the arguments of this
paper the design condition is needed for the components in every case, and
for the prediction rate of Theorem~\ref{thm:main} only when $s \ge 1/2$.}
% E5f (E1.15; D54): a frase acima qualifica a condicao de desenho; a
% introducao (D50) nao muda.
Two further statements complete the picture.
The first dispenses with the design condition altogether, at the price of a
slower rate \textcolor{colR1}{when $s \ge 1/2$}, and is the unconditional statement of this paper.

% derivations/05-taxas.tex, Proposicao 4 (numeracao global). Substituida em
% k = 4 pelo Corolario 14 de derivations/08-blocos.tex (D47(c)); a do lasso,
% emendada por D49(a), e a Proposition S8.1.
% E1.14 (integrada pela E5e, 2026-10-04): os itens (iii) e (iv) do Corolario 14,
% s' <= 1/4 com s > 1/2, s' <= s/2 com s < 1/2 e a fronteira s = 1/2, estao na
% Proposition S6.1 do supp (prop:slow-ceiling); aqui, uma frase depois da
% proposicao, e a frase de D49(b) revista (D52).
{\color{gray}\removedk
\begin{propositionold}[Slow rate, no design condition]
Let Assumptions~\ref{ass:sample}, \ref{ass:basis} and \ref{ass:besov} hold,
with $\Amat$ of rank $p$ and $\penn$ as in the first display of
Section~\ref{sec:main}. Then
$\normn{\widehat{f}-f}^{2} = O_p\{\penn\norm{\thetastar}_{1} +
\mathcal{B}_{J_n}^{2} + p/n\}$, where
$\norm{\thetastar}_{1} \le pqC_g\sum_{\lev<J_n}2^{\lev(1/2-s)}$ does not
depend on $\pi$. In particular, if $s > 1/2$ then
$\normn{\widehat{f}-f}^{2} = O_p\{(J_n/n)^{1/2}\}$ for every
$J_n \ge \log_2 n/(4\seff)$; and if $s < 1/2$, the choice
$2^{J_n} \asymp (n/J_n)^{1/(1-2s+4\seff)}$ gives
$\normn{\widehat{f}-f}^{2} = O_p\{(J_n/n)^{2\seff/(1-2s+4\seff)}\}$, an
exponent that equals $2s/(2s+1)$ when $\pi \ge 2$.
\end{propositionold}
}
% derivations/08-blocos.tex, Corolario 14 (numeracao global); D47(c), D49(b).
{\color{colR1}
\begin{proposition}[Slow rate, no design condition]
\label{prop:slow}
Let Assumptions~\ref{ass:sample}, \ref{ass:basis} and \ref{ass:besov} hold,
together with the upper bounds in Assumptions~\ref{ass:X} and \ref{ass:U}
($\max_{\cv}|X_{\cv}| \le B_X$, $\lmaxx\{\E(\bX\bX\T\mid\bU)\} \le \kappa_2$ and
$p_{\bU} \le C_U$), with $\Amat$ of rank $p$, the weights as in
Theorem~\ref{thm:main}, $\pen = \lamG$ as in \eqref{eq:lambda} with any constant
$\Lambda > \kappa_2C_U$, $2^{J_n} \le 2n$ and
$pq2^{J_n}\log(pq2^{J_n})/n \to 0$. Then
\[
  \normn{\widehat{f}-f}^{2} = O_p\bigl\{\lamG\nGw{\thetastar} +
  \mathcal{B}_{J_n}^{2} + p/n\bigr\},
\]
with $\lamG\nGw{\thetastar} \le \bar\varrho\,\lamGone\nGone{\thetastar}
\le \bar\varrho\,\lamGone\norm{\thetastar}_{1}$, where $\nGone{\thetastar} = \sum_{G}\norm{\thetastar_{G}}_{2}$ and
$(\lamGone)^{2}$ is of order $\sigma^{2}\log n/n$. Write
$\vartheta_\pi = \min(1-1/\pi,\,1/2)$.
\begin{enumerate}
\item[(i)] If $s > 1/2$ and $\seff > 1/4$, then $\nGone{\thetastar} = O(1)$
and $\normn{\widehat{f}-f}^{2} = O_p\{(\log n/n)^{1/2}\}$ for every
$J_n \ge \log_2 n/(4\seff)$ that satisfies the conditions above, and one
does.
\item[(ii)] If $s < 1/2$ and $\seff > s/2$, then
$\nGone{\thetastar} = O\{1 + \bn^{-\vartheta_\pi}2^{J_n(1/2-s)}\}$, and the
choice $2^{J_n} \asymp \{n/\bn^{(2/\pi-1)_{+}}\}^{1/(1-2s+4\seff)}$ gives
\[
  \normn{\widehat{f}-f}^{2}
  = O_p\bigl\{n^{-2\seff/(1-2s+4\seff)}
  (\log n)^{(2/\pi-1)_{+}2\seff/(1-2s+4\seff)}\bigr\} ,
\]
which, when $\pi \ge 2$, is $n^{-2s/(2s+1)}$, the rate of
Theorems~\ref{thm:main} and \ref{thm:sieve}, with no logarithmic factor.
\end{enumerate}
The proof uses no lower bound on $\Gramhat$, and no property of $\Gramhat$
other than \textcolor{gray}{\sout{$\lmaxx(\Gramhat) \le \Lambda$}} an upper bound on
$\lmaxx(\Gramhat)$.
\end{proposition}
}

\textcolor{colR1}{The bound on $\nGone{\thetastar}$ is Lemma~S6.2 in the
Supplementary Material. Through the chain to $\norm{\thetastar}_{1}$,
Proposition~\ref{prop:slow} reproduces the slow rate of the lasso
(Proposition~S8.1), with the constant $\bar\varrho$: a chunk of up to $2\bn-1$
coefficients enters $\lamGone$ at the order of the union bound over the chunks.
What the chain throws away is the count: at the fine levels
$\nGone{\thetastar}$ is smaller than $\norm{\thetastar}_{1}$ by the factor
$\bn^{-\vartheta_\pi}$, which in (ii) cancels the factor $\bn^{1/2}$ that the
size of the chunks puts in $\lamGone$, entirely when $\pi \ge 2$ and in part
when $\pi < 2$; in (i) the coarse chunk dominates both norms and nothing
changes. No lower bound on $\Gramhat$ enters: the only property of $\Gramhat$
that the proof uses is \textcolor{gray}{\sout{$\lmaxx(\Gramhat) \le \Lambda$}} an upper bound on
$\lmaxx(\Gramhat)$, here $\lmaxx(\Gramhat) \le \Lambda$, which makes the
deterministic level \eqref{eq:lambda} admissible and holds with probability
tending to one under the upper bounds of Assumptions~\ref{ass:X}
and~\ref{ass:U} alone, so that neither $c_U$ nor $\kappa_1$ is needed. The
conditions $\seff > 1/4$ in (i) and $\seff > s/2$ in (ii) bite only when
$\pi < 2$, and make the resolution of the balance admissible.
% E5e (D52): os pares que (i) e (ii) excluem e a fronteira s = 1/2 vao ao supp.
The pairs that these conditions leave out, $\seff \le 1/4$ with $s > 1/2$ and
$\seff \le s/2$ with $s < 1/2$, all with $\pi < 2$, and the boundary $s = 1/2$
call for a resolution above $n/\log n$, where the column norms and
$\lmaxx(\Gramhat)$ grow with $2^{J_n}/n$; Proposition~S6.1 in the
Supplementary Material covers them with the level inflated to
$\lamGp = (\mu^{\Gcal}_{J_n})^{1/2}\lamG$, where
$\mu^{\Gcal}_{J} = 1 + c_{\mu}2^{J}\log(2D)/n$ for a constant $c_{\mu}$ given
there, so that the first display of Proposition~\ref{prop:slow}, with $\lamGp$
in place of $\lamG$, holds for every $\seff > 0$. There, too, the only
property of $\Gramhat$ that the proof uses is an upper bound,
$\lmaxx(\Gramhat) \le \Lambda\mu^{\Gcal}_{J_n}$, which holds with probability
tending to one with no condition on $J_n$ and, again, without $c_U$ or
$\kappa_1$.
% E5f (E1.15; D54): o Corolario 15 de derivations/08-blocos.tex, que e a
% Proposition S6.2 do supp; a versao do lasso (Proposicao 8 do 05) e a
% Proposition S8.2.
When $s < 1/2$ the slow rate improves further. Read at the comparator that
keeps the chunks with $\norm{\thetastar_{G}}_{2} > \pen w_{G}$ and sets the
others to zero, Theorem~\ref{thm:oracle}(i) charges the penalty only for the
chunks kept and leaves the energy of the others to the bias, where an upper
bound on $\lmaxx(\Gram)$ controls it; with $\seff > s/(2s+1)$ and $J_n$ in the
window \eqref{eq:window}, this gives the prediction rate of
Theorem~\ref{thm:main} with no design condition (Proposition~S6.2 in the
Supplementary Material). When $s \ge 1/2$ no comparator does better than the
order $\lamGone \asymp (\log n/n)^{1/2}$ of part (i): a coarse coefficient of
fixed size costs its product with $\lamGone$ in the penalty or its square in
the bias. Neither reading reaches the components, which need the design
condition (Corollary~\ref{cor:components}).}

The second is the approximation lemma that produces Theorem~\ref{thm:main}.
\textcolor{gray}{\sout{For a vector $\btheta$ let $\theta_{(1)} \ge \theta_{(2)} \ge \cdots$ be the
nonincreasing rearrangement of the absolute values of its entries and let
$\btheta^{(M)}$ retain the $M$ largest of them.}}
\textcolor{colR1}{For $\btheta \in \R^{d}$ and $\eta > 0$, let
$\Rid(\btheta;\eta) = \sum_{G\in\Gcal}\min(\norm{\btheta_{G}}_{2}^{2},\eta)$
be the ideal risk over chunks: the risk of an oracle that keeps a chunk when
its energy exceeds $\eta$, at the price $\eta$, and drops it otherwise.}

% derivations/05-taxas.tex, Lema 9 (numeracao global). Substituido em k = 4 pelo
% Lema 15 de derivations/08-blocos.tex; o do lasso e o Lemma S8.1.
{\color{gray}\removedk
\begin{lemmaold}[Besov regularity implies compressibility]
Under Assumptions~\ref{ass:basis} and \ref{ass:besov} with
$\pi(s+1/2) > 1$, the oracle satisfies, for every $J$ and every
$k \ge 1$,
\[
  \theta^{*}_{(k)} \;\le\; C_\tau\,k^{-1/\tau},
  \qquad \tau = \frac{1}{s+1/2},
  \qquad C_\tau = C_g(pq\,A_{s,\pi})^{s+1/2},
  \qquad A_{s,\pi} = 1 + \frac{1}{1-2^{1-\pi(s+1/2)}},
\]
with constants that do not depend on $J$ and with $\tau \in (0,2)$ for every
$s > 0$. Consequently
$\sum_{k>M}\theta^{*2}_{(k)} \le \tau(2-\tau)^{-1}C_\tau^{2}M^{1-2/\tau}$ for
every $M \ge 1$.
\end{lemmaold}
}
% derivations/08-blocos.tex, Lema 15 (numeracao global).
{\color{colR1}
\begin{lemma}[Besov regularity bounds the ideal risk over chunks]
\label{lem:weak}
Under Assumptions~\ref{ass:basis} and \ref{ass:besov}, let $b \ge 1$ and let
$\Gcal$ be a partition in which, in each block, the levels with $2^{\lev} < b$
form at most one chunk and every level with $2^{\lev} \ge b$ is cut into
chunks contained in it with at least $b$ coefficients each, as the partition of
Section~\ref{sec:estimator} is. Then, for every $J$ and every $\eta > 0$,
\[
  \Rid(\thetastar;\eta) \;\le\; pq\Bigl\{\AG\,C_g^{\tau}
  \Bigl(\frac{\eta}{b}\Bigr)^{2s/(2s+1)}b^{(2/\pi-1)_{+}/(2s+1)} + \eta\Bigr\},
  \qquad \tau = \frac{1}{s+1/2},
\]
with $\AG = 2 + \{1-2^{1-\pi(s+1/2)}\}^{-1}$ if $\pi \le 2$ and
$\AG = 2 + (1-2^{-2s})^{-1}$ if $\pi \ge 2$.
\end{lemma}
}

\textcolor{gray}{\sout{The inclusion of Besov balls in weak $\ell_\tau$ balls is classical
\mbox{\citep[Proposition~9.15]{johnstone2019gaussian}}; what the lemma adds is the
constant $C_\tau$, uniform in $J$, for the oracle of the product design, which
is the form in which the proof of Theorem~\mbox{\ref{thm:main}} uses it.}}
\textcolor{colR1}{The exponent is the one of Lemma~4 in the supplementary material
of \citet{Klopp-Pensky-2015}, but the assumption is weaker: their class is a
weighted $\ell_\nu$ ball over all the coefficients, which implies
Assumption~\ref{ass:besov} level by level, and Assumption~\ref{ass:besov} does
not need the sum over the levels. The coarse chunk costs at most one $\eta$
per block.}

\textcolor{gray}{\sout{Theorem~\mbox{\ref{thm:main}} follows by reading Theorem~\mbox{\ref{thm:oracle}}(ii) not at
the oracle but at the best $M$ term approximation of it, which is legitimate
because the proof of that theorem uses the comparator only through the
decomposition of $\bY$ and the size of its support. Balancing the two terms
that result, $\pen^{2}M/\gtil$ from estimation and
$M^{1-2/\tau}$ from the truncation, gives the effective dimension}}
{\color{gray}\removedk
\[
  \Mstar_n \;\asymp\; \Bigl(\frac{n}{J_n}\Bigr)^{\tau/2}
  \;\ll\; d_n \asymp pq\,2^{J_n},
\]
}
\textcolor{gray}{\sout{and a bound of order $\penn^{2-\tau}$, which is the rate of
Theorem~\mbox{\ref{thm:main}}. The three inequalities that define the window
\mbox{\eqref{eq:window}} are, in order, that the minimum in $M$ be interior, that the
approximation bias not dominate, and that the requirement of
Proposition~\mbox{\ref{prop:gram}}(ii) hold. Nothing in this argument says that the
estimator selects $\Mstar_n$ coordinates: no minimal signal or
irrepresentability condition is used, and what compressibility buys is a rate
governed by $\Mstar_n$ rather than by $d_n$, not the recovery of a support.}}
\textcolor{colR1}{Theorem~\ref{thm:main} follows by reading
Theorem~\ref{thm:oracle}(iii) not at the oracle but at the comparator that
keeps the chunks of $\thetastar$ whose squared norm exceeds a level $\eta$ and
sets the others to zero, which the theorem allows because it holds for every
comparator at once. The estimation term then counts $\eta$ for each chunk
kept, and the truncation leaves the energy of the others, so that the bound is
of the order of $\Rid(\thetastar;\eta)$ (Corollary~S5.1 in the Supplementary
Material), with $\eta$ of the order of $\max_{G}(\pen w_{G})^{2}$, that is of
$\bn/n$. Lemma~\ref{lem:weak} with $b = \bn$ turns this into
$n^{-2s/(2s+1)}\bn^{(2/\pi-1)_{+}/(2s+1)}$: the chunk pays $\bn$ in the level
$\eta$ and returns it through $(\eta/b)^{2s/(2s+1)}$, so that the cost is of
order $1/n$ per coefficient, against $\log n/n$ for the lasso, and the
logarithm of the union bound over the columns is absorbed by the size of the
chunks, which is why they must have at least of the order of $\log n$
coefficients. When $\pi < 2$ the bound charges
$\bn^{(2/\pi-1)/(2s+1)}$ to sequences whose large coefficients sit one per
chunk; a signal spread over a level pays nothing for the chunks. The two
inequalities that define the window \eqref{eq:window} are that the
approximation bias not dominate and that the requirement of
Proposition~\ref{prop:gram}(ii) hold. Nothing in this argument says that the
estimator selects the chunks kept by the comparator: no minimal signal or
irrepresentability condition is used.}
\textcolor{colR1}{The recovery that matters for the model is not of coordinates
but of blocks, and it is taken up next.}

% derivations/06-selecao-limiar.tex, Lema 13, Hipotese S e Corolario 8
% (numeracao global); D32. O Lema 13 e a prova estao no supp (Secao S7).
{\color{colR1}
\subsection{Which modulators act on which coefficients}
\label{sec:selection}

The structure of the model is the set of blocks whose component does not
vanish,
\[
  \mathcal{S} \;=\; \bigl\{(\cv,\md) :\; \nu_{\cv\md} > 0\bigr\},
  \qquad
  \nu_{\cv\md} = \bigl\lVert \gcomp{\cv}{\md} \bigr\rVert_{L_2[0,1]},
\]
and to know it is to know which modulators act on which coefficients. The
\textcolor{gray}{\sout{lasso}} estimator \eqref{eq:estimator} is not claimed to recover it: its
penalty removes chunks and not blocks, under Assumption~\ref{ass:besov} the
oracle is dense, and nothing prevents the fit from keeping nonzero
coefficients in blocks whose component is zero. What the preceding results do give is that every block is
estimated at the rate of Corollary~\ref{cor:components}, and that is enough to
recover $\mathcal{S}$ in a second step, by a threshold on the norms of the fitted
components. The procedure is estimation followed by a threshold. By the
orthonormality of the basis, the norm of a fitted component is the Euclidean
norm of its block of coefficients,
\[
  \widehat{\nu}_{\cv\md}
  = \bigl\lVert \widehat{g}_{\cv\md} \bigr\rVert_{L_2[0,1]}
  = \norm{\thetahat_{\cv\md,\cdot}}_{2},
  \qquad
  \widehat{\mathcal{S}}(t) = \bigl\{(\cv,\md) :\; \widehat{\nu}_{\cv\md} > t\bigr\},
\]
so that the statistic costs nothing once the estimator has been fit. Write
\textcolor{gray}{\sout{$\rho_n = (\log n/n)^{\seff/(2\seff+1)}$}} $\rhoG = n^{-\seff/(2\seff+1)}$ for the square root of the rate of
Theorem~\ref{thm:sieve}.

\begin{assumption}[Block separation]
\label{ass:sep}
The components may depend on $n$, within the class of
Assumption~\ref{ass:besov} and with its constants fixed, and so may $\mathcal{S}$; the
smallest active block, $\nu_{\min,n} = \min_{(\cv,\md)\in \mathcal{S}} \nu_{\cv\md}$, with
$\nu_{\min,n} = \infty$ when $\mathcal{S}$ is empty, satisfies \textcolor{gray}{\sout{$\nu_{\min,n}/\rho_n \to \infty$}} $\nu_{\min,n}/\rhoG \to \infty$.
\end{assumption}

\begin{corollary}[Block structure by thresholding]
\label{cor:threshold}
Let the assumptions of Theorem~\ref{thm:sieve} hold, with $J_n$ and \textcolor{gray}{\sout{$\penn$}} $\pen = \lamG$ as
there.
\begin{enumerate}
\item[(i)] (Screening.) If the thresholds satisfy \textcolor{gray}{\sout{$t_n/\rho_n \to \infty$}} $t_n/\rhoG \to \infty$,
then $\Prob\{\widehat{\mathcal{S}}(t_n) \subseteq \mathcal{S}\} \to 1$.
\item[(ii)] (Recovery.) If, in addition, Assumption~\ref{ass:sep} holds and
$\limsup_n t_n/\nu_{\min,n} < 1$, then $\Prob\{\widehat{\mathcal{S}}(t_n) = \mathcal{S}\} \to 1$.
Such thresholds always exist; one is \textcolor{gray}{\sout{$t_n = (\rho_n\nu_{\min,n})^{1/2}$}} $t_n = (\rhoG\nu_{\min,n})^{1/2}$.
\item[(iii)] \textcolor{colR1}{(A window at fixed $n$.) With $\alpha$ the level in
\eqref{eq:lambda}, for every $n$ and every $\alpha' \in (0,1)$ there is an
event, of probability at least
$1 - \alpha - \alpha' - \omega_n$ with $\omega_n \to 0$ not depending on the
components, on which $\widehat{\mathcal{S}}(t) = \mathcal{S}$ for every $t$
with \textcolor{gray}{\sout{$\bar{\Delta}_n \le t < \nu_{\min,n} - \bar{\Delta}_n$}} $\bar{\Delta}^{\Gcal}_n \le t < \nu_{\min,n} - \bar{\Delta}^{\Gcal}_n$; here
\textcolor{gray}{\sout{$\bar{\Delta}_n \le c(\alpha,\alpha')\rho_n$}} $\bar{\Delta}^{\Gcal}_n \le c(\alpha,\alpha')\rhoG$ is an explicit bound on
$\max_{(\cv,\md)}\lVert \widehat{g}_{\cv\md} - \gcomp{\cv}{\md}
\rVert_{L_2[0,1]}$, given in the proof (Section~S7 of the Supplementary
Material).}
\end{enumerate}
\end{corollary}

\textcolor{colR1}{Parts (i) and (ii) hold also at the resolution of
Theorem~\ref{thm:main}, with \textcolor{gray}{\sout{$\rho_n$}} $\rhoG$ replaced by
\textcolor{gray}{\sout{$\tilde\rho_n = (\log n/n)^{s/(2s+1)}$}} $\rhoGt = n^{-s/(2s+1)}(\log n)^{(2/\pi-1)_{+}/\{2(2s+1)\}}$, the square root of the rate of that
theorem, both in the condition on $t_n$ and in Assumption~\ref{ass:sep}.
Since \textcolor{gray}{\sout{$\tilde\rho_n \le \rho_n$}} $\rhoGt \le \rhoG$ for large $n$, with a gap in the exponent when $\pi < 2$,
the separation required is weaker exactly where the components are spatially
inhomogeneous. Part (iii) is stated at the resolution of
Theorem~\ref{thm:sieve}, where its bound is explicit. The componentwise rate of
the lasso at its own resolution is $(\log n/n)^{\seff/(2\seff+1)}$
(Corollary~S8.1 in the Supplementary Material), so that the separation
required here is weaker than for the lasso by the factor
$(\log n)^{\seff/(2\seff+1)}$.}

Part (i) uses no separation: removing the null blocks is free, and
Assumption~\ref{ass:sep} is spent entirely on not losing a small active block.
Under a fixed target $\nu_{\min,n}$ is a positive constant and the assumption holds
automatically, because \textcolor{gray}{\sout{$\rho_n \to 0$}} $\rhoG \to 0$; it carries content in a triangular
array, where it says how fast the smallest active block may shrink. It bears
on the components and not on $(\bX,\bU)$, so that no design condition beyond
Proposition~\ref{prop:gram} and no irrepresentability condition enters; a
separation condition on the components plays the same role in the selection
of variables in nonparametric additive models \citep{Huang-Horowitz-Wei-2010}.
The proof rests on a deterministic sandwich (Lemma~S7.1 in the Supplementary
Material): if every component is estimated within $\bar{\Delta}$ in $L_2[0,1]$, then
$\{\nu_{\cv\md} > t + \bar{\Delta}\} \subseteq \widehat{\mathcal{S}}(t) \subseteq
\{\nu_{\cv\md} > t - \bar{\Delta}\}$, so that every $t \in [\bar{\Delta}, \nu_{\min,n} - \bar{\Delta})$ recovers $\mathcal{S}$,
and that window cannot be widened by the argument. Two things are left out.
The window involves $\sigma$, the design constant and $\seff$, so the
corollary does not say how to choose $t_n$ from the data \textcolor{colR1}{(Section~\ref{sec:threshold}
uses the one standard error rule of cross-validation, for which no theorem is
given here)}; and it says nothing
about which wavelets are active inside a block.

% derivations/06-selecao-limiar.tex, Lema 16 e Corolario 13 (numeracao global;
% adendo de E1.12); D45. O Lema 16 e a prova estao no supp (Secao S7).
Thresholding is also the estimator used for prediction
(Section~\ref{sec:threshold}). Write $\thetahat^{(t)}$ for $\thetahat$ with the
blocks outside $\widehat{\mathcal{S}}(t)$ set to zero, and
$\widehat{g}^{(t)}_{\cv\md}$ and $\widehat{f}^{(t)} = \Amat\chat +
\Bmat\thetahat^{(t)}$ for the resulting components and fit, so that the levels
and the blocks kept are those of the fit. Setting to zero an active block
whose fitted norm is at most $t$ costs at most its norm, which is at most
$t + \lVert\widehat{g}_{\cv\md} - \gcomp{\cv}{\md}\rVert_{L_2[0,1]}$, setting to
zero a null block only reduces the error, and the blocks kept do not change
(Lemma~S7.2 in the Supplementary Material). With the rates of the fit this
gives the risk of the thresholded estimator.

\begin{corollary}[Risk of the thresholded estimator]
\label{cor:thrisk}
Let $(\chat,\thetahat)$ be either the fit of Theorem~\ref{thm:sieve}, with
$r_n = \rhoG$, or the lasso at the resolution of Theorem~S8.2 in the
Supplementary Material, with $r_n = (\log n/n)^{\seff/(2\seff+1)}$, and let
$t_n \ge 0$.
\begin{enumerate}
\item[(i)] With no further assumption,
$\sum_{\cv,\md}\lVert\widehat{g}^{(t_n)}_{\cv\md}-\gcomp{\cv}{\md}\rVert_{L_2[0,1]}^{2}
= O_p(r_n^{2} + t_n^{2})$ and
$\normn{\widehat{f}^{(t_n)}-f}^{2} = O_p(r_n^{2} + t_n^{2})$.
\item[(ii)] If Assumption~\ref{ass:sep} holds with $r_n$ in place of $\rhoG$
and $\limsup_n t_n/\nu_{\min,n} < 1$, then with probability tending to one
$\mathcal{S} \subseteq \widehat{\mathcal{S}}(t_n)$,
\[
  \sum_{\cv,\md}\lVert\widehat{g}^{(t_n)}_{\cv\md}-\gcomp{\cv}{\md}\rVert_{L_2[0,1]}^{2}
  \;\le\; \sum_{\cv,\md}\lVert\widehat{g}_{\cv\md}-\gcomp{\cv}{\md}\rVert_{L_2[0,1]}^{2}
\]
and
\[
  \normn{\widehat{f}^{(t_n)}-f}^{2} \;\le\;
  \frac{4\Lambda}{\gamma}\normn{\widehat{f}-f_{J_n}}^{2} + 2\normn{\bb}^{2},
\]
so that the thresholded estimator keeps the rate $r_n^{2}$ of the fit. If
moreover $t_n/r_n \to \infty$, then
$\widehat{\mathcal{S}}(t_n) = \mathcal{S}$ with probability tending to one,
and the null blocks no longer contribute to the error of the components.
\end{enumerate}
\end{corollary}

Under a fixed target Assumption~\ref{ass:sep} holds automatically, and (ii)
holds for every $t_n \to 0$. The factor $4\Lambda/\gamma$ in (ii) bounds the
condition number of $\Gramhat$, through which the bound passes from
prediction to coefficients and back; it is a price of the argument, and the
corollary says what any threshold costs, not which threshold a data driven
rule selects.
}

%%%%%%%%%%%%%%%%%%%%%%%%%%%%%%%%%%%%%%%%%%%%%%%%%%%%%%%%%%%%%%%%%%%%%%%%%%%%%%
\section{Computation and Tuning}
\label{sec:computation}

\subsection{The design matrix and the solver}
\label{sec:solver}

The design \eqref{eq:design} is built one block at a time: the $n \times \NJ$
matrix of the wavelets of $U_{\md}$ is evaluated once per modulator, by the
cascade algorithm of \citet{Daubechies-Lagarias-1991} \textcolor{colR1}{as implemented in
the \texttt{WaveBased} package \citep{Montoril-WaveBased-2026}}, and then multiplied
rowwise by $X_{\cv}$. Since a compactly supported wavelet of level $\lev$ is
nonzero on an interval of length $(L-1)2^{-\lev}$, at most $L-1$ of the
$2^{\lev}$ columns of a level are nonzero at a given observation, so a block
has at most $n(L-1)J$ nonzero entries out of $n\NJ$ and the design is stored in
sparse form as soon as $J$ is moderate. Evaluating the basis once on a grid and
looking the values up removes the cascade from the inner loop of a
simulation.

Problem \eqref{eq:estimator} is then a \textcolor{gray}{\sout{lasso}}\textcolor{colR1}{group lasso} with $p$ unpenalized coordinates,
% E5f (L10; D54): as duas citacoes do grpreg dentro dos grupos azuis de k = 4;
% a nota [VERIFICAR] de E5d saiu sem cinza, por ser nota de trabalho.
which \textcolor{gray}{\sout{coordinate descent}}\textcolor{colR1}{group descent \citep{Breheny-Huang-2015}} solves along a path of penalty levels; we use the
\textcolor{gray}{\sout{\texttt{glmnet}}}\textcolor{colR1}{\texttt{grpreg}} package\textcolor{gray}{\sout{ \mbox{\citep{friedman2010regularization}}}} \textcolor{colR1}{\citep{Breheny-Miller-grpreg-2026}} \textcolor{colR1}{of R
\citep{RCoreTeam-2026}}. Two points of scale deserve care, because they are
what ties the software to the theory. First, \textcolor{colR1}{the solver standardizes
the columns and orthonormalizes each chunk before penalizing it
% L11 (D55, pergunta 49(b)): Simon & Tibshirani (2012), eq. (1.4) e Section 2.
\citep{Simon-Tibshirani-2012}, so that what
it computes is \eqref{eq:estimator} with $\norm{\btheta_{G}}_{2}$ replaced by
$\norm{Q_{G}^{1/2}\btheta_{G}}_{2}$, where $Q_{G}$ is the Gram matrix of the
columns of the chunk after centering, which is the estimator of Remark~S5.2 in
the Supplementary Material; and its objective,
$(2n)^{-1}\norm{\bY - \Amat\bc - \Bmat\btheta}_{2}^{2} +
\pen\sum_{G}w_{G}\norm{Q_{G}^{1/2}\btheta_{G}}_{2}$, is half that of
\eqref{eq:estimator}, so that the penalty level it reports is the $\pen$ of
\eqref{eq:estimator}, with no conversion. Second, the solver always fits an
intercept of its own, which carries the level of the constant covariate when
there is one. For the lasso option, solved by coordinate descent with the
\texttt{glmnet} package \citep{friedman2010regularization},} the penalty level reported for a
fit is the $\pen$ of \textcolor{gray}{\sout{\mbox{\eqref{eq:estimator}}}}\textcolor{colR1}{the lasso problem} and not the one of the solver, which
rescales its own by the ratio between the total number of columns and the
number of penalized ones when some coordinates carry zero penalty\textcolor{gray}{\sout{; the
conversion is checked by verifying the Karush-Kuhn-Tucker conditions of
\mbox{\eqref{eq:estimator}} at every point of the path. Second,}}\textcolor{colR1}{, and} when one linear
covariate is the constant, its column has zero variance, and the solver may
either drop it or share the level with its own intercept; the fit is the same
either way, but the level $\lvl{\cv}$ of the constant covariate is only
recovered after the intercept is added back to it. \textcolor{colR1}{For both
penalties the scale is checked by verifying the Karush-Kuhn-Tucker conditions
of the problem at every point of the path.}

\subsection{Rescaling and the boundary}
\label{sec:boundary}

The theory places $\bU$ in $[0,1]^q$ by assumption. In practice each modulator
is mapped to that interval by a monotone transformation\textcolor{gray}{\sout{, and,}}\textcolor{colR1}{.
By default the image is $[0,1]$ itself, which is the setting of
Section~\ref{sec:theory}. Optionally,} as in the
wavelet estimators of \citet*{Montoril-Chang-Vidakovic-2019}, the image is
$[\epsilon, 1-\epsilon]$ with $\epsilon$ \textcolor{gray}{\sout{of the order $2^{-J}$}}\textcolor{colR1}{fixed}, which keeps the
observations away from the \textcolor{gray}{\sout{strip}}\textcolor{colR1}{point} where the periodization of
\eqref{eq:periodized} \textcolor{gray}{\sout{leaves an artifact}}\textcolor{colR1}{joins the two ends of the interval}; the price is that the fitted
component integrates to zero over the rescaled range and not over $[0,1]$, a
shift of level and not of shape. \textcolor{colR1}{The margin is a finite sample
device, and the reason to allow it is approximation: with the data away from
the ends, a component needs to be regular only on its support, where it can be
replaced by a periodic extension that matches at $0 \equiv 1$, built from a
bounded extension operator \citep[Chapter~3]{Triebel-1983}. This removes the
stall of Proposition~S3.1 and restores the bias of Lemma~\ref{lem:bias}, with
a Besov constant of order $\epsilon^{-(s-1/\pi)}$, an order that, by the
embedding of Besov into H\"older spaces, cannot be improved when the component
takes different values at the two ends of its support. Two features keep the
margin out of the theory. It has to be fixed, since with $\epsilon$ of the
order $2^{-J}$ the exponent guaranteed for the bias is no better than that of
periodization without a margin; and it degrades the design, since the
wavelets supported inside the margin are invisible to the data, and the
smallest eigenvalue of the Gram matrix of the basis restricted to the support
of the data vanishes as soon as $2^{J} \ge (L-1)/(2\epsilon)$, so that
Proposition~\ref{prop:gram} does not hold as stated.} The alternative is the basis adapted to the
interval of \citet*{cohen1993wavelets}, which reproduces polynomials on $[0,1]$
and characterizes Besov regularity without any periodicity
\citep[Section 3.9]{Cohen-2003}, so that Lemma~\ref{lem:bias} holds with the
Besov space of the interval and the same proof. Its cost is structural: the
coarsest level must be large enough for the two boundary zones to be disjoint,
the constant is no longer a single column but a combination of the scaling
functions, and the constraint of \eqref{eq:additive} has to be imposed by
reparametrization, which leaves $pq(2^{\jzero}-1)$ unpenalized parameters
instead of none. We therefore keep the periodized basis as the default and
offer the interval basis as an option.

\subsection{Choosing the resolution and the penalty}
\label{sec:tuning}

The pair $(J,\pen)$ is chosen jointly. The default is $K$\textcolor{colR1}{-}fold \textcolor{gray}{\sout{cross
validation}}\textcolor{colR1}{cross-validation} on a grid of resolutions, with the folds drawn once and reused for
every $J$, one design and one penalty path per resolution, and the usual two
readings of the curve: the minimizer of the \textcolor{gray}{\sout{cross validated}}\textcolor{colR1}{cross-validated} error and the
largest penalty within one standard error of it\textcolor{colR1}{; the threshold
of Section~\ref{sec:threshold} is applied to the fit at the minimizer. In the
software the whole procedure, from the cross-validation over $(J,\pen)$ to the
threshold, is the single call \texttt{cv.wafc(x, u, y)}}. Two cheaper rules are also
available \textcolor{colR1}{for the lasso option}, the Bayesian information criterion and its extended version, with
degrees of freedom equal to the number of nonzero coefficients plus the $p$
levels, which is the usual unbiased count for the lasso.

The theory supplies a fourth rule, since \eqref{eq:Jn} and \eqref{eq:lambda}
are explicit. It is not a practical competitor as it stands, because it needs
$\sigma$, which has to be estimated, and $\seff$, which is unknown, and
because it is a statement about order\textcolor{gray}{\sout{, so that at moderate sample sizes it can
sit above the resolution that minimizes the realized error, the pre-asymptotic
regime being wide when $\seff$ is small}}\textcolor{colR1}{. In our experiments} \textcolor{colR1}{with the
lasso option} \textcolor{colR1}{the penalty}
\textcolor{gray}{\sout{\mbox{\eqref{eq:lambda}}}} \textcolor{colR1}{of its theory} \textcolor{colR1}{is 7 to 13 times the one that minimizes the realized error,
and it is the penalty, more than the resolution, that makes the rule costly;
the resolution}
\textcolor{gray}{\sout{\mbox{\eqref{eq:Jn}}}} \textcolor{colR1}{of its theory} \textcolor{colR1}{misses the best one by up to two levels, above
or below according to $\seff$}.
% E5d: os numeros acima sao do lasso (E2.3, E1.6); com o bloco como padrao,
% passaram a dizer isso. Uma medida equivalente para o block lasso e da E5b.
% E5b: ancorar os numeros acima numa tabela da Secao 5 ou do suplemento
% (E2.3: lambda_n de 7 a 13 vezes o otimo; no nao homogeneo, J_n sozinho
% custa 2% e 18%, lambda_n sozinho 30% e 54%; J_n dois niveis abaixo com
% s' = 3/2, dois acima com s' = 1/4 em E1.6).
Its interest is as a reference for
what the \textcolor{gray}{\sout{cross validated}}\textcolor{colR1}{cross-validated} choice ought to approach\textcolor{gray}{\sout{, and the difference between
the two is a finite sample quantity worth measuring}}.\textcolor{gray}{\sout{ The four rules are
compared in the simulation study.}}
% E5i (D68(a)): a frase das quatro regras saiu, porque o estudo da Secao 5 (D60)
% roda so a validacao cruzada do WAFC; o comentario E5b logo acima segue com a E5b.
% E5j (D72(a)): a oracao "worth measuring" saiu pelo mesmo motivo.
% E5j (D72(b)): o "$K$-fold" do 4.3 marca so o hifen em colR1; com "fold" tachado e
% "-fold" inserido, a linha passava 17 pt da margem na versao marcada.

A remark on the \textcolor{gray}{\sout{structural variant}}\textcolor{colR1}{other penalties}. \textcolor{colR1}{The lasso option, with
$\norm{\btheta}_1$ in place of $\nGw{\btheta}$, is fit by \texttt{glmnet}; its
theory is in Section~S8 of the Supplementary Material, where its rates carry
an extra logarithmic factor, and the threshold of Section~\ref{sec:threshold}
applies to it unchanged.} Replacing \textcolor{gray}{\sout{$\norm{\btheta}_1$}}\textcolor{colR1}{$\nGw{\btheta}$} in
\eqref{eq:estimator} by the sparse group lasso penalty of
\citet{Simon-Friedman-Hastie-Tibshirani-2013}, with one group per block
$(\cv,\md)$ \textcolor{colR1}{and fit with the \texttt{sparsegl} package
\citep{Liang-Cohen-SolonHeinsfeld-Pestilli-McDonald-2024}}, adds a second tuning parameter and removes entire components
rather than isolated \textcolor{gray}{\sout{coefficients}}\textcolor{colR1}{chunks}, which is what a reader interested in
knowing which modulators act on which coefficients would want\textcolor{colR1}{, although for}
\textcolor{gray}{\sout{the lasso itself}} \textcolor{colR1}{\eqref{eq:estimator}} \textcolor{colR1}{that question is already answered by thresholding the block
norms (Corollary~\ref{cor:threshold})}. It is a
different estimator, with a different oracle inequality
\citep{Lounici-Pontil-vandeGeer-Tsybakov-2011}, and the theory above does not
cover it; the choice between the two is settled empirically.

% derivations/06-selecao-limiar.tex, adendo de E1.12; D45. A regra e a
% wafc_threshold(rule = "cv1se") de wafc/R/threshold.R, padrao do cv.wafc() (D48).
{\color{colR1}
\subsection{Choosing the threshold}
\label{sec:threshold}

The threshold is chosen on the folds of Section~\ref{sec:tuning}, at the pair
$(J,\pen)$ already selected: each fold is refitted at that pair, with no new
search, thresholded at every candidate $t$ and scored on the observations it
left out. A thresholded fit is constant in $t$ between two consecutive block
norms, so the candidates are exhaustive and not a grid. The default is the
largest $t$ whose cross-validated error is within one standard error, between
folds, of the smallest, the one standard error rule of the penalty carried to
the threshold; the minimizer of the cross-validated error is available as an
option aimed at prediction. The fit is not refitted after the threshold: the
levels and the blocks kept are those of the fit, which is the estimator of
Corollaries~\ref{cor:threshold} and~\ref{cor:thrisk}. Neither rule has a
theorem here: Corollary~\ref{cor:threshold} gives the window of thresholds
that recovers the structure, and Corollary~\ref{cor:thrisk} what any threshold
costs.
}

%%%%%%%%%%%%%%%%%%%%%%%%%%%%%%%%%%%%%%%%%%%%%%%%%%%%%%%%%%%%%%%%%%%%%%%%%%%%%%
% E5h (k = 5): a Secao 5 sem os resultados. O desenho e o de D60, D62 e D63
% (docs/plano-projeto.md, E4.2), como o compendio wafc-studies/ o executa
% (config/study.yaml, core.yaml, arms.yaml, scale.yaml; R/data.R, R/methods.R,
% R/metrics.R; o gerador e wafc/R/dgp.R); os metodos e as grades sao os de D34,
% D41, D44 a D46, D56 e D63. O VCBART roda no compendio e nao entra no texto
% (D60(d)). [E4.3] espera o piloto (o modelo maior; a grade de n = 2000 e a de D69,
% aplicada na E5i);
% [E4.4], a producao. A frase do terceiro braco e a da E4.1b (ESTADO.md), com o
% pareamento escrito para p - 1 <= q, que cobre as duas celulas; a das
% covariaveis gaussianas e a de D62.
{\color{colR1}
\section{Simulation}
\label{sec:simulation}

\subsection{Design}
\label{sec:sim-design}

The samples are drawn from model \eqref{eq:model} with $X_1 \equiv 1$, the
other linear covariates independent and standard normal, the modulators
independent of $\bX$, independent of each other and uniform on $[0,1]$, and
Gaussian errors. The levels are $(\lvl{1},\lvl{2},\lvl{3},\lvl{4}) =
(1, 2, -1.5, 0.5)$, repeated cyclically when $p > 4$, and each nonzero
component is a test function centered and scaled to have unit norm in
$L_2[0,1]$. Unless it is fixed, the standard deviation of the error is that of
the regression function over the sample divided by a signal to noise ratio.

The core design has five cells. In the smooth cell, with $(p,q) = (3,2)$, the
first coefficient is additive in the two modulators, with $\gcomp{1}{1}$ and
$\gcomp{1}{2}$ proportional to $\sin(2\pi u)$ and to the cubic
$u^{3} - 1.4u^{2} + 0.4u$, the second varies with $U_1$, with $\gcomp{2}{1}$
proportional to $\cos(4\pi u)$, the other components are zero, and the signal
to noise ratio is 4. The uneven cell keeps this structure, the cosine and the
ratio, and replaces the sine by a Gaussian bump of standard deviation $0.10$
at $0.40$ minus $0.8$ times one of standard deviation $0.018$ at $0.72$, and
the cubic by the chirp $\sin\{2\pi(u + 3.5u^{2})\}$, whose frequency goes from
one to eight cycles over the interval; both are multiplied by an infinitely
differentiable window that equals one on $[0.08, 0.92]$ and vanishes with all
its derivatives at $0$ and $1$ (Section~S9 of the Supplementary Material).
These components are infinitely differentiable, within the smoothness
conditions of the spline theory of \citet{Xue-Yang-2006}, but the scale of two
of them changes along the interval. The inhomogeneous cell keeps the
structure, with the bumps, blocks and heavisine functions of
\citet{Donoho-Johnstone-1994} in place of the sine, the cubic and the cosine,
and signal to noise ratio 3. The mixed cell places the components of the
inhomogeneous one in a model with $(p,q) = (4,4)$, so that 3 of its 16 blocks
are active, with the same ratio. In the null cell every component is zero and
the model is a linear regression with $(p,q) = (3,2)$; the standard deviation
of the error is fixed at $0.62$, about a quarter of that of the regression
function.

In the terms of Assumption~\ref{ass:besov}, the effective regularity is
$\seff = 3/2$ in the smooth cell, where the cubic has different slopes at the
two ends of the interval and its periodic extension has a corner, and
$\seff = 1/2$ in the inhomogeneous and mixed cells, where blocks and heavisine
jump inside the interval. The components of the uneven cell have infinitely
differentiable periodic extensions, and their regularity is limited only by
the vanishing moments of the wavelet.

Each cell is run at $n = 250$, $500$, $1000$ and $2000$, with 100 replicates
of each pair of cell and sample size. A replicate consists of a training
sample of size $n$, a test sample of size 2000 from the same law, with errors
at the standard deviation of the training sample, and a partition of the
training sample into 10 folds; the three are shared by every method, so that
the comparisons within a replicate are paired.

Three arms change one factor at a time in the inhomogeneous and the mixed
cells, at $n = 1000$ and with 100 replicates each. In the first, the signal to
noise ratio is 1 instead of 3. In the second, the modulators are coupled by a
Gaussian copula with correlation $0.5$ between every pair, which keeps the
uniform marginals and makes the cross terms of the Gram matrix between
distinct modulators nonzero; the density of this copula is unbounded at two
corners of the unit square and vanishes at the other two, so that this arm
lies outside Assumption~\ref{ass:U}. In the third arm the linear covariates
depend on the modulating ones: each non-constant covariate is drawn as
$X_\cv = (1-\rho^2)^{1/2} \zeta_\cv + \rho\, h(U_{\cv-1})$ with $\rho = 0.5$,
where $\zeta_\cv$ is standard normal and $h(u) = \sqrt{12}\,(u - 1/2)$ is the
standard score of a uniform modulating covariate, so that $X_2$ is paired with
$U_1$, the modulating covariate of its own coefficient. Each $X_\cv$ keeps
unit variance and has correlation $\rho$ with its pair, and given $\bU$ the
non-constant covariates have covariance $(1-\rho^2)I$, so that the eigenvalue
condition of Assumption~\ref{ass:X} holds while the Gram matrix of the
products no longer factorizes as a Kronecker product. Gaussian covariates are
not bounded, in this arm as in every cell, so that the design lies outside the
bound $B_X$ of Assumption~\ref{ass:X}; what the third arm exercises is the
eigenvalue condition, the part of that assumption that allows the covariates
to depend on the modulators. A replicate of an arm is drawn from the seeds of
the same replicate of its core cell, and the first and the third arms share
with it the modulators, the draws $\zeta_\cv$ and the standardized errors.

A larger model [E4.3: run if the pilot measures its cost below about ten hours
of processor time] has $(p,q) = (6,4)$ at $n = 1000$, with 24 blocks of which
6 are active: the active structure of the inhomogeneous cell twice, on
$(X_1,X_2)\times(U_1,U_2)$ and on $(X_3,X_4)\times(U_3,U_4)$, with $\beta_5$
and $\beta_6$ constant and signal to noise ratio 3.

Every sample is derived from a single master seed, through an injective
encoding of the cell, the sample size, the replicate and the stream. The seeds
are independent of those of the pilot runs on which the estimator, its tuning
rules and the grids below up to $n = 1000$ were settled, so that the study does not evaluate
the method on the samples on which it was chosen.

\subsection{Methods and tuning}
\label{sec:sim-methods}

Every method is tuned in its own terms, by the criteria it is used with, so
that a difference between two methods is not a difference of tuning. The WAFC
is the estimator \eqref{eq:estimator} with $b = \lceil\log n\rceil$ and
weights $|G|^{1/2}$, on the periodized extremal phase wavelets of
\citet{Daubechies-1992} with four vanishing moments ($L = 8$), evaluated by
table and with no margin; $(J,\pen)$ is chosen by 10-fold cross-validation
(Section~\ref{sec:tuning}) and the threshold by the one standard error rule on
the same folds (Section~\ref{sec:threshold}), the default of the software. The
row WAFC (cv) comes from the same search, with the threshold that minimizes
the cross-validated error, the option aimed at prediction. Seven other rows
enter the tables, five competitors and two references.

GAM (REML) and GAM (GCV) fit the additive varying coefficient model by
penalized regression splines with the \texttt{mgcv} package
\citep{Wood-2017,Wood-mgcv-2025}, through its function \texttt{bam}
\citep{Wood-Goude-Shaw-2015}, with one smooth of $U_\md$ multiplied by $X_\cv$
for every block and an extra penalty on the null space of each smooth, which
lets a smooth be shrunk to zero \citep{Marra-Wood-2011}. The first chooses the
smoothing parameters by restricted maximum likelihood and the second by
generalized cross-validation, and each chooses the basis dimension $k$, common
to the smooths, from a grid by its own criterion, which is the full search of
\citet{Ruppert-2002} and the search of \citet[Section~2.3]{Pya-Wood-2016}. GAM
(GCV) is not run in the mixed cell or in the larger model, where a single fit
at the top of its grid takes more than an hour.

The B-spline group lasso expands each component in cubic B-splines, centered
on the training sample, and fits a group lasso \citep{Yuan-Lin-2006} with one
group per block and the levels unpenalized, with \texttt{grpreg}; the number
of basis functions per block, $2^J$, and the penalty are chosen by
cross-validation on the folds. It differs from the WAFC in the basis and in
the unit of the penalty, a whole block instead of a chunk.

The adaptive spline gives each block a cubic spline with its own knots, chosen
from the data, in the spirit of \citet{Wang-Jiang-Liu-2024}, who choose the
knots by dynamic programming in a model with a single modulator. Here the
knots of a block, at most 12 among 40 candidates at quantiles of its
modulator, are inserted one at a time by the Bayesian information criterion,
inside a backfitting cycle over the blocks, so that this row stands in for
their method and does not implement it.

The block lasso (KP) is the estimator of \citet{Klopp-Pensky-2015} on the
wavelet design of the WAFC, with their chunk size $\lceil\log n\rceil$ and the
levels penalized as one more group, as in their norm. It departs from their
estimator in the chunks, which stay inside a resolution level, each coarse
level being a chunk of its own, and in the weights, which are the $|G|^{1/2}$
of the solver; it differs from the WAFC in the penalized levels and in the
cut of the chunks. Its $(J,\pen)$ is chosen by cross-validation, with no
threshold.

The two references are the oracle support, the estimator
\eqref{eq:estimator} on the columns of the active blocks only, with $(J,\pen)$
chosen by cross-validation and no threshold, whose distance to the WAFC is the
price of not knowing the structure; and the constant coefficients, least
squares on $\bX$ alone, which is the correct model in the null cell, where the
oracle support reduces to it, and a floor in the others. The coordinatewise
lasso of Section~\ref{sec:computation}, at the pair chosen by
cross-validation and thresholded by the one standard error rule, is reported
in Section~S9 of the Supplementary Material.

The grids were fixed by the pilot runs and do not depend on the cell:
$J \in \{2,\dots,8\}$ for the methods on the wavelet design, the same $J$ for
the $2^J$ B-splines per block, and $k \in \{5, 10, 20, 40, 80\}$ for the two
searches of \texttt{mgcv}, one grid for both criteria, which avoids powers of
two so that the splines are not tuned on the grid of the wavelets. A
resolution is capped by the number of distinct points of the modulator: at
$n = 250$ the 255 wavelets of $J = 8$ exceed the 249 distinct points that a
modulator has on the circle after rescaling, where its two extremes meet at
$0 \equiv 1$, so the level is built as $J = 7$ and the effective grid there
is $\{2,\dots,7\}$, and the dimension 256 of the B-splines is dropped. At
$n = 2000$, and only there, the grids are extended before the study, to
$J \in \{2,\dots,9\}$ for the methods on the wavelet design and for the
B-splines, and with $k = 120$ added to the grid of \texttt{mgcv}, because the
resolution that the theory prescribes grows with $n$, as in \eqref{eq:Jn}.

\subsection{Measures}
\label{sec:sim-measures}

Every method is asked the same questions on the test sample and on a grid of
256 equally spaced points in the observed range of each modulator. The
prediction error is the root mean squared difference between the fitted and
the true regression function on the test sample. The error of the components
is the integrated squared error summed over the blocks, with the fit and the
truth both centered on the grid, because every method fixes the level of a
component by its own convention; it is split between the active blocks and
the null ones, where it measures what a method puts where there is nothing.
The two errors are reported separately, since a method can predict well with
wrong components, or the reverse. The structure is read from the blocks a
method keeps, as the frequency $\Prob(\widehat{\mathcal{S}} = \mathcal{S})$ of
the replicates in which they are exactly the active ones, beside the mean
numbers of null blocks kept and of active blocks dropped; a smooth of
\texttt{mgcv} counts as kept when its effective degrees of freedom exceed
$0.1$, since the extra penalty shrinks it toward zero but never to zero
exactly. The time is the elapsed time of a fit, the search over the grids and
the threshold included, in a single process [E4.4: the processor]. Means over
the replicates are reported with their standard errors; since the methods
share the samples of a replicate, the Supplementary Material also gives the
median, within replicate, of the ratio of each error to that of the WAFC, and
the fraction of replicates in which a method has the smaller error.

\subsection{Results}
\label{sec:sim-results}

Table~\ref{tab:sim-main} gives the errors and the recovery of the structure at
$n = 1000$, and Figure~\ref{fig:sim-components} the fitted components of the
inhomogeneous cell. [E4.4: the reading of Table~\ref{tab:sim-main} and
Figure~\ref{fig:sim-components}.] Section~S9 of the Supplementary Material
gives the errors at every sample size, the arms [E4.3: and the larger model],
the structure, the resolution and the dimension chosen with the frequency of
choices at the top of a grid, the time, and the slope of the logarithm of the
errors against $\log n$. [E7: the address of the compendium that reproduces
the study.]

% tables/sim-main.tex -- Table 1 do ms_5.tex (E5h, 2026-10-06; D21: cada tabela
% num \input proprio, para que mandar ao supp seja mover uma linha).
% So o lugar: o conteudo vem da E4.4 (results/tables/, do summary.csv e do
% paired.csv do compendio wafc-studies/). Proposta de layout, a decidir na E5b:
% linhas = as nove linhas do corpo, com os rotulos de config/study.yaml (WAFC,
% WAFC (cv), GAM (REML), GAM (GCV), B-spline group lasso, Adaptive spline,
% Block lasso (KP), Oracle support, Constant coefficients); colunas = as cinco
% celulas do nucleo; tres paineis empilhados em n = 1000, (a) erro de predicao,
% (b) ISE das componentes, media (erro-padrao), (c) P(S-hat = S). O GAM (GCV)
% fica sem numero na mixed (nao roda ali). Sem paisagem (instrucoes da SS).
\begin{table}[t]
\caption{Simulation at $n = 1000$, 100 replicates per cell. [E4.4: root mean
squared prediction error and integrated squared error of the components, mean
and standard error, by method and cell, and the frequency with which the
blocks kept are exactly the active ones.]}
\label{tab:sim-main}
\centering
\fbox{\parbox{0.9\textwidth}{\centering\small [E4.4: the table, from
\texttt{results/tables/}.]}}
\end{table}

% figures/sim-components.tex -- Figure 1 do ms_5.tex (E5h, 2026-10-06; D21).
% So o lugar: a figura vem da E4.4, das curvas que o compendio guarda
% (config/study.yaml, `curves`: wafc.cv1se e gam.reml em n = 1000). O plano
% (E4.2) pede g-hat sobre a verdade em bumps e blocks, WAFC contra GAM (REML);
% qual replica ou que resumo pontual se desenha e decisao da E4.4/E5b. Figura
% legivel em preto e branco, sem referencia a cor na legenda (SS).
% E5k (D75(d)): \color{colR1} dentro do float, que sai preto no grupo de fora.
\begin{figure}[t]
\color{colR1}
\centering
\fbox{\parbox[c][5cm][c]{0.9\textwidth}{\centering\small [E4.4: the figure,
from the curves of the compendium.]}}
\caption{Inhomogeneous cell, $n = 1000$: the components $\gcomp{1}{1}$
(bumps, left) and $\gcomp{1}{2}$ (blocks, right), the truth and the fits of
the WAFC and of GAM (REML). [E4.4: the line types, the replicate or the
pointwise summary over the 100 replicates that is drawn, and the finding.]}
\label{fig:sim-components}
\end{figure}

}

%%%%%%%%%%%%%%%%%%%%%%%%%%%%%%%%%%%%%%%%%%%%%%%%%%%%%%%%%%%%%%%%%%%%%%%%%%%%%%
% E5h (k = 5): a Secao 6 sem os numeros. Os dados, o modelo e as leituras sao os
% de D56, D57, D59, D64, D66 e D67, como o compendio os executa
% (wafc-studies/config/application.yaml, R/application_data.R, R/application.R,
% R/application_report.R, data/README.md); as fontes sao as de L12 e L13. Os
% numeros ficam como [E6.2: ...], a trocar pelos de readings-, prediction- e
% structure-marylebone.ukair.csv da rodada. E5j (2026-10-07): os de marylebone
% trocados pelos de results/tables/ (handoff-E5j.md, cada numero e a sua
% coluna); a frase dos buracos no 6.1 (D73, "spacing" no lugar de "support"),
% o pico na segunda metade de 2003 e a serie de Carslaw ate dezembro de 2003
% (respostas do autor na E5j). Os da beijing.heat (6.4) ficam. A frase de Opsomer, Wang e Yang (2001)
% nao entra (D65). As atribuicoes que as licencas pedem estao nos
% Acknowledgements.
{\color{colR1}
\section{Application}
\label{sec:application}

\subsection{Data and model}
\label{sec:app-data}

At a roadside site the total oxidant $\mathrm{OX} = \mathrm{NO_2} +
\mathrm{O_3}$ is close to linear in the concentration of the nitrogen oxides
NOx \citep[Section~2.3]{Clapp-Jenkin-2001}: the intercept is the regional
background of oxidant, and the slope is the local contribution to it, which
\citet[Section~2.4]{Carslaw-Beevers-2005} use as an estimate of the fraction
of NOx emitted directly as $\mathrm{NO_2}$, the primary fraction. The slope
is an estimate and not the fraction itself: it also carries the
$\mathrm{NO_2}$ formed by the thermal reaction of NO with $\mathrm{O_2}$ and
from emitted nitrous acid, and it is lower at night
\citep[Section~2.3]{Clapp-Jenkin-2001}. At Marylebone Road, London,
\citet[Sections~3.1 and 3.2, Figure~3(a)]{Carslaw-2005} estimated the primary
fraction at about 10\% by volume from 1997 to 2002, rising through 2002 and
2003 to about 23\% at the end of 2003, and related the rise mainly, though not
wholly, to the particle filters fitted to the buses (his Sections~3.3
and~3.4). A coefficient of NOx that varies with the date can show such a
change without a model for its timing, and a second modulator asks whether
the coefficient also moves with the weather.

The data are the hourly concentrations of NOx, expressed as $\mathrm{NO_2}$,
of $\mathrm{NO_2}$ and of $\mathrm{O_3}$ at Marylebone Road, site MY1 of the
Automatic Urban and Rural Network, from January 1998 to June 2005
\citep{Defra-UKAIR-2026}, converted to parts per billion (ppb) with the
conversion factors of the same department \citep{Defra-2005}, and the hourly wind speed at 10 m of the ERA5
reanalysis \citep{Hersbach-2020,Hersbach-ERA5-2023} at the nearest grid point,
retrieved through the Open-Meteo service \citep{Zippenfenig-OpenMeteo-2024}
and interpolated to the middle of each hour. The 61,280 hours with the four
variables observed are used. The concentrations are published under the Open
Government Licence and the wind under the Creative Commons Attribution 4.0
licence; the attributions they require are given in the Acknowledgements.

The response is OX in ppb, the linear covariates are $X_1 \equiv 1$ and
$X_2$, NOx in units of 100 ppb, so that $\beta_2$ reads as a percentage of
NOx, and the modulators are the date and the wind speed, each acting on both
coefficients, so that $p = q = 2$ and the model has four blocks. Both
modulators are mapped to $[0,1]$ by their range. The date is not periodic, and
the periodized basis joins the end of the series to its start
(Section~\ref{sec:boundary}), so that the fits are read away from the first
and the last 60 days. The series also has intervals without measurements,
the longest of 43 days in the summer of 1998; over those longer than the
spacing of the wavelets of the finest level chosen, about 11 days, the
estimate is not determined by the data and is not shown.

\subsection{Fitting and evaluation}
\label{sec:app-fit}

The errors are correlated over hours and days, and a test sample of hours
scattered among the training hours would reward a method for learning the
level of each week. The unit of every split is therefore a week. The
prediction error is estimated over 20 random partitions that hold out 30\% of
the weeks, and inside each training sample the weeks are spread over 10
folds, which every method tuned by cross-validation uses; the figure shows
the fits to the whole sample, with folds formed in the same way. The
partitions are drawn from the master seed of the simulation study and are
independent of the exploratory partitions on which this data set was chosen.
Each method is tuned in its own terms, as in Section~\ref{sec:sim-methods}:
the WAFC chooses $(J,\pen)$, with $J \in \{2,\dots,8\}$, and the threshold by
cross-validation on these folds, \texttt{mgcv} chooses the smoothing
parameters and $k \in \{5,10,20,40,80\}$ by REML or by GCV, and the constant
coefficients have nothing to tune. None of the methods models the dependence
of the errors. The error of a method is the root mean squared error on the
held-out weeks, averaged over the partitions, with the standard error of
\citet{Nadeau-Bengio-2003}, which corrects for the overlap between the
training samples of different partitions.

\subsection{The step of 2003}
\label{sec:app-step}

Figure~\ref{fig:marylebone} shows each coefficient along each modulator, with
the component of the other modulator at its sample mean. The background of
oxidant follows the annual cycle, over a range of 23 to 24~ppb in the fits to
the whole sample. The slope on NOx is flat until mid 2002 and then rises, by
8.9 percentage points from its level before 2002 to its level from 2004 on in
the fit to the whole sample, and by 9.3 to 10.3 points over the 20
partitions; it has completed a quarter of the rise by January 2003 (0.16 to
0.36 over the partitions) and 90\% of it by the end of July 2003, in the fit
to the whole sample and in every partition. The two thresholds keep the same
blocks along the date and give the same rise, and the spline fitted by REML
gives a rise of 10.3 points, 90\% of it by the beginning of August 2003. At
this site \citet[Figure~3(a)]{Carslaw-2005} estimated a rise of the primary
fraction from about 10\% to about 23\% by volume over 2002 and 2003; the step
of the fit is measured from 2004 on, past the peak described next, and the
rise to that peak is 13.4 points (11.7 to 13.8 over the partitions). Two
departures from a single step appear in every fit of the WAFC and of the
spline fitted by REML, a dip at the beginning of 2002 and a peak in the
second half of 2003: the dip also appears in the monthly estimates of
\citet[Figure~3(a)]{Carslaw-2005}, whose series ends in December 2003,
before the slope falls back from the peak. Neither is interpreted here.

The wind speed modulates the slope by 0.8 to 1.8 percentage points, 9\% to
17\% of the step, in the fits to the whole sample that keep that block. The
WAFC drops that block, in the fit to the whole sample and in every partition,
while the threshold of the smallest cross-validated error keeps it in 19 of
the 20 partitions and the splines fitted by REML and by GCV in 18 and 19, so
the fit does not support the absence of that modulation, only its size,
which is small beside the step. Both modulators act on the background in
every fit of the threshold of the smallest cross-validated error and of the
two splines, the wind speed over a range of 5.7 to 7.8~ppb in the fits to
the whole sample; the WAFC keeps the date and drops the wind speed in the
fit to the whole sample, and keeps the wind speed in 12 of the 20
partitions. On the held-out weeks, the WAFC with the threshold of the
smallest cross-validated error predicts as well as the better of the two
splines: its root mean squared error is 11.52~ppb against 11.47 for the
spline fitted by REML, which has the smaller mean error of the two; the
ratio to the better spline of each partition is 1.005 on average, and the
paired difference, 0.05~ppb with standard error 0.11, is negative in 8 of
the 20 partitions. The WAFC is behind by 2.5\%, with an error of 11.75~ppb
and a paired difference of 0.28~ppb (standard error 0.22). Section~S10 of the
Supplementary Material gives the errors of every method, the frequency with
which each block is kept over the partitions, and the choices of $J$, $\pen$,
the threshold and $k$.

% figures/app-marylebone.tex -- Figure 2 do ms_5.tex (E5h, 2026-10-06; D21).
% A figura e a marylebone.ukair.pdf do compendio (R/application_report.R,
% split00 e as 20 particoes; D67(d): beta_l remontado, 60 dias cortados em cada
% ponta da data). A legenda e o rascunho da E6.2 no ESTADO.md, com tres mudancas:
% os numeros ganharam o prefixo "E6.2:" e "in every partition" foi aos colchetes
% (as particoes do compendio sao novas, D67(b)); "as does the rise from about 10
% to about 23 percent by volume that Carslaw estimated" passou a uma frase
% propria, porque comparava uma subida com outra pela sintaxe do verbo; e "Fig."
% passou a "Figure", como no corpo. Sem referencia a cor (SS).
% A legenda inteira do rascunho, com a comparacao com Carslaw (2005), o mergulho
% e o vento, nao cabe numa pagina em espaco duplo (Float too large por 220 pt
% com a caixa de 7 cm; a curta ainda passa 15 pt com 7 cm, e a caixa ficou em
% 5 cm); essas frases ficam no texto da Secao 6.3, e o trecho
% tirado daqui fica abaixo para a E5b, se a legenda voltar a ser longa:
%   grey: the range of the WAFC over 20 fits on 70\% of the weeks. The
%   background follows the annual cycle. The slope, an estimate of the fraction
%   of NOx emitted directly as $\mathrm{NO_2}$
%   \citep{Clapp-Jenkin-2001,Carslaw-Beevers-2005}, rises by [E6.2: 9.7]
%   percentage points from 2002 to 2004, [E6.2: a quarter] of the rise by January
%   2003 and [E6.2: 90\%] by July 2003 [E6.2: in every partition]. At this site,
%   \citet[Figure~3(a)]{Carslaw-2005} estimated a rise from about 10 to about 23
%   percent by volume over 2002 and 2003. The dip at the beginning of 2002 also
%   appears in his monthly estimates, whose series ends before the peak at the end
%   of 2003. The wind speed modulates the slope by about [E6.2: 17\%] of the step;
%   the WAFC drops that block, and the threshold of the smallest cross-validated
%   error and the spline keep it. 
\begin{figure}[tp]
\centering
\fbox{\parbox[c][5cm][c]{0.9\textwidth}{\centering\small [E6.2: the figure
\texttt{marylebone.ukair.pdf} of the compendium.]}}
\caption{Marylebone Road, London, hourly, 1998--2005: total oxidant
$\mathrm{NO_2} + \mathrm{O_3}$ regressed on NOx with coefficients modulated by
the date and the wind speed. Each panel shows $\widehat\beta_\ell$ along one
modulating covariate, with the component of the other at its sample mean:
the oxidant background (top, ppb) and the slope on NOx (bottom, percent of
NOx), along the date (left) and the wind speed (right). Solid: the WAFC
fitted to the whole sample; dashed: the WAFC with the threshold of the
smallest cross-validated error; dot-dashed: \texttt{mgcv} with REML;
grey: the range of the WAFC over 20 fits on 70\% of the weeks. The
background follows the annual cycle, and the slope, an estimate of the
fraction of NOx emitted directly as $\mathrm{NO_2}$, rises by [E6.2: 9.7]
percentage points from 2002 to 2004. The first and the last 60 days are not shown:
there the periodized basis joins the end of the series to its start.}
\label{fig:marylebone}
\end{figure}


\subsection{A second application}
\label{sec:app-beijing}

Section~S10 of the Supplementary Material applies the method to the hourly
PM2.5 at the site Dongsi in Beijing, from March 2013 to February 2017, in the
Beijing Multi-Site Air Quality data of the UCI Machine Learning Repository
\citep{Chen-UCI-2017, Zhang-Guo-Dong-He-Xu-Chen-2017}:
the logarithm of PM2.5 is regressed on the wind speed, the temperature and
the pressure, with coefficients modulated by the day of the year, a periodic
modulator, and the relative humidity, over 34,287 hours. The question there
is prediction. The coefficients change with the season, around the nominal
dates of the heating season, 15 November to 15 March
\citep[Section~7]{Liang-Zou-Guo-Li-Zhang-Zhang-Huang-Chen-2015}, and on the
held-out weeks the WAFC [E6.2: the prediction against the better of the two
splines].
}

%%%%%%%%%%%%%%%%%%%%%%%%%%%%%%%%%%%%%%%%%%%%%%%%%%%%%%%%%%%%%%%%%%%%%%%%%%%%%%
\section*{Supplementary Materials}

The online Supplementary Material contains the proofs of all the results
stated in the paper, together with the auxiliary lemmas they use: the
identification argument and the properties of the periodized basis, the
approximation lemmas in $L_2$ and in the supremum norm and the proposition on
the cost of periodization, the population and empirical bounds for the Gram
matrix of the product design, the \textcolor{gray}{\sout{four lemmas}}\textcolor{colR1}{properties of the
partition into chunks, the calibration of the penalty over chunks and the
other lemmas} behind the oracle inequality \textcolor{colR1}{and its form for the
estimator that the software computes},
\textcolor{gray}{\sout{and}} the \textcolor{gray}{\sout{two lemmas}}\textcolor{colR1}{lemmas} behind the rates\textcolor{colR1}{, \textcolor{gray}{\sout{and}} the \textcolor{gray}{\sout{lemma}} two lemmas behind
the recovery of the block structure by thresholding} \textcolor{colR1}{and the risk of
the thresholded estimator, and the theory of the coordinatewise lasso, with
its proofs}\textcolor{colR1}{. It also gives the test functions and the larger
model of the simulation study, with its results at every sample size and in
every arm, and further results of the application, with the second one, to
data from Beijing}.

%%%%%%%%%%%%%%%%%%%%%%%%%%%%%%%%%%%%%%%%%%%%%%%%%%%%%%%%%%%%%%%%%%%%%%%%%%%%%%
\section*{Acknowledgements}

Write the acknowledgements here.
\par
% E5h (k = 5): as atribuicoes que as licencas dos dados pedem, copiadas como os
% licenciadores as dao (L12 para o UK-AIR, o ERA5 e o Open-Meteo; a da UCI do
% wafc-studies/data/README.md); por isso "licenced" fica na grafia britanica.
{\color{colR1}
The data of Section~\ref{sec:application} are used under the following
attributions. \copyright{} Crown 2026 copyright Defra via
uk-air.defra.gov.uk, licenced under the Open Government Licence (OGL).
Generated using or contains modified Copernicus Climate Change Service
information 2026. Neither the European Commission nor ECMWF is responsible
for any use that may be made of the Copernicus information or data it
contains. Hourly ERA5 wind data were retrieved through the Open-Meteo
historical weather API
\citetext{\citealp{Zippenfenig-OpenMeteo-2024}; \url{https://open-meteo.com/}},
licensed under CC BY 4.0 (\url{https://creativecommons.org/licenses/by/4.0/});
the instantaneous values were linearly interpolated to the middle of each
hour. Beijing Multi-Site Air Quality data, UCI Machine Learning Repository
(\url{https://doi.org/10.24432/C5RK5G}), licensed under CC BY 4.0
(\url{https://creativecommons.org/licenses/by/4.0/}); only the site Dongsi is
used.
\par
}

%%%%%%%%%%%%%%%%%%%%%%%%%%%%%%%%%%%%%%%%%%%%%%%%%%%%%%%%%%%%%%%%%%%%%%%%%%%%%%

\bibhang=1.7pc
\bibsep=2pt
\fontsize{9}{14pt plus.8pt minus .6pt}\selectfont
\renewcommand\bibname{\large \bf References}
\expandafter\ifx\csname
natexlab\endcsname\relax\def\natexlab#1{#1}\fi
\expandafter\ifx\csname url\endcsname\relax
  \def\url#1{\texttt{#1}}\fi
\expandafter\ifx\csname urlprefix\endcsname\relax\def\urlprefix{URL}\fi

\bibliographystyle{chicago}
\bibliography{references_5}

%-------------------------------------------
\vskip .65cm
\noindent
first author affiliation
\vskip 2pt
\noindent
E-mail: (first author email)
\vskip 2pt

\noindent
second author affiliation
\vskip 2pt
\noindent
E-mail: (second author email)

\end{document}
